\documentclass[openany]{book}
\usepackage{amsmath, amsthm, amssymb}
\usepackage{hyperref} % 关键！启用目录跳转
\usepackage{geometry} % 添加geometry包控制页边距
\usepackage{indentfirst} %首行缩进
\usepackage{tikz-cd}
\usepackage{mathtools}

\geometry{
    a4paper,           % 纸张大小
    left=20mm,         % 左边距
    right=20mm,        % 右边距
    top=25mm,          % 上边距
    bottom=25mm,       % 下边距
    bindingoffset=0mm  % 装订边距设为0，使奇偶页对称
}
\hypersetup{colorlinks=true, linkcolor=blue}
\newtheorem{theorem}{Theorem}[chapter] % 定理，按章编号，如 Theorem 1.1, 1.2
\newtheorem{lemma}[theorem]{Lemma}     % 引理，与定理共用计数器 Lemma 1.3
\newtheorem{proposition}[theorem]{Proposition} % 命题
\newtheorem{corollary}[theorem]{Corollary}   % 推论

% 定义、例子、备注通常使用不同的字体样式
\theoremstyle{definition} % 将后续环境设置为“定义”样式（正文字体）
\newtheorem{definition}[theorem]{Definition} % 定义
\newtheorem{example}[theorem]{Example}       % 例子
\newtheorem{exercise}[theorem]{Exercise}     % 练习

\theoremstyle{remark}
% 让remark单独编号，不与其他定理环境共享计数器
\newtheorem{remark}{Remark}[chapter] % 备注
\newtheorem{note}[remark]{Note} % 注释

\title{D-modules}
\begin{document}
\maketitle
\tableofcontents
\chapter{Basics of D-modules on smooth affine algebraic varieties}
\chapter{Characteristic varieties$\&$cycles of D-modules}
\chapter{Filtered rings and modules}
\chapter{Some symplectic geometry of affine $T^{\ast}X$}
\chapter{Gabber’s Theorem: Involutivity of Characteristic varieties}
\chapter{Holonomicity and regular holonomicity}
\section{Holonomic $\mathcal{D}$-modules}
Let $X = \text{spec} R$ be a smooth affine scheme with algebraic coordinates $(x_1, \ldots, x_n)$ and let $M$ be a finite generated $\mathcal{D}_X$-module. We already know $\text{ch}(M) \subset T^*X$ is a conic involutive (or coisotropic) subvariety. In particular, $\dim \text{ch}(M) \geq n$. 

\begin{definition}
    $M$ is called $\textbf{holonomic}$ if $\dim \text{ch}(M) = n$.
\end{definition}

\begin{corollary}
    If $M$ is a holonomic $\mathcal{D}_X$-module, then $\text{ch}(M) \subset T^*X$ is a Lagrangian subvariety, and $$CC(M) = \sum_{E_i \subset X\atop \text{irr subvarieties}} n_i T^*_{E_i} X$$
\end{corollary}

\begin{remark}
    Recall that Lagrangian means $W=W^{\perp}\iff \mathrm{dim}_{k}W=n$, and we define $V\subseteq T^{\ast}X$ is Lagrangian if for any smooth closed point $p\in V$, we have $\mathcal{T}_{V,p}\subseteq \mathcal{T}_{T^{\ast}X,p}$ is Lagrangian. Here $T^{\ast}_{E_{i}}X$ is the conormal bundle of irreducible subvariety $E_{i}$. 
\end{remark}

\begin{example}
    \quad
    \begin{enumerate}
        \item $R$ is a $\mathcal{D}_X$-module. $CC(R) = T^*_X X\simeq X \Rightarrow R$ is holonomic.
        \item Let $M$ be an integrable connection on $X$, we obtain that $ch(M)=T^{\ast}_{X}X$ and $\mathrm{dim}_{k}V(\xi_{1},\dots,\xi_{n})=n$, we have $M$ is holonomic. 
    \end{enumerate}
\end{example}

\section{The de Rham Complexes and Analytification}

Let $M$ be a $D_X$-module, and denote $\partial_i = \partial_{x_i}$ be the vector field along $x_i$.  
Define a complex $$K(M; \partial_1) = [M \xrightarrow{\partial_1} M].$$Here $[M\xrightarrow{\partial_{1}} M]$ means a complex concentrated on $\mathrm{deg}=0$ and $\mathrm{deg}=1$, with $\partial_{1}$ as coboundary map. 

\quad

\begin{note}[The cone of a morphism of complexes]\label{note 6.2}
Let $f: A^\bullet \longrightarrow B^\bullet$ be a morphism of complexes:

\begin{center}
 \begin{tikzcd}
 \cdots \arrow[r] & A^{n-1} \arrow[r, "d_A^{n-1}"] \arrow[d, "f^{n-1}"] & A^n \arrow[r, "d_A^n"] \arrow[d, "f^n"] & A^{n+1} \arrow[r] \arrow[d, "f^{n+1}"] & \cdots \\
 \cdots \arrow[r] & B^{n-1} \arrow[r, "d_B^{n-1}"] & B^n \arrow[r, "d_B^n"] & B^{n+1} \arrow[r] & \cdots
 \end{tikzcd}
\end{center}

then $\mathrm{Cone}(f)=C^{\bullet}$ defined by $C^n = A^n \oplus B^{n-1}$ with differential $$d_C^n = \begin{pmatrix}
 d_A^n & 0 \\
 f^n & -d_B^{n-1}
 \end{pmatrix}: A^n \oplus B^{n-1} \longrightarrow A^{n+1} \oplus B^n.$$The cone construction gives a short exact sequence of complexes:
\begin{center}
\begin{tikzcd}
0 \arrow[r] & B^\bullet[-1] \arrow[r] & Cone(f)^\bullet \arrow[r] & A^\bullet \arrow[r] & 0.
\end{tikzcd}
\end{center}
\end{note}


With Note \ref{note 6.2}, for $k \geq 2$, we inductively define     $$K(M; \partial_1, \cdots, \partial_k) = \text{cone}\left(K(M; \partial_1, \cdots, \partial_{k-1}) \xrightarrow{\partial_k} K(M; \partial_1, \cdots, \partial_{k-1})\right).$$ Then the $\textbf{de Rham complex}$ of $M$ is   $$K(M; \partial_1, \cdots, \partial_n),$$ denoted as $\mathrm{DR}(M).$

\begin{remark}
    In fact, it's Koszul complex with $(\partial_{1},\dots,\partial_{n})$ a finite sequence of central elements in $M$, defined in a mapping cone way.   
\end{remark}

Meanwhile, we have a coordinate-free way to define $\mathrm{DR}(M)$: $$M \to \Omega^{1}_{X} \otimes_{R} M \to \Omega^{2}_{X} \otimes_{R} M \to \cdots \to \Omega^{n}_{X} \otimes_{R} M.$$With $m\rightarrow \sum\limits_{i}dx_{i}\otimes \partial_{i}(m)$, and $\eta \otimes m \mapsto \sum dx_i \wedge \eta \otimes \partial_i(m)$ and so on...

\begin{remark}
    Here $\sum\limits_{i} dx_i \otimes \partial_i \in \Omega^1 \otimes_R \mathcal{T}_X$ is canonical i.e. independent of coordinates.
\end{remark}

\begin{lemma}
    We have $$\mathrm{DR}(M) \cong K(M; \partial_1, \ldots, \partial_n).$$The coordinate-free de Rham complex is isomorphic to the Koszul complex on the partial derivatives.
\end{lemma}
\begin{proof}
    When $n = 2$,  $\mathrm{DR}(M):0 \to M \to \Omega^{1}_{X} \otimes M \to \Omega^{2}_{X} \otimes M \to 0.$ $$K(M; \partial_{1}, \partial_{2})= M \to M \oplus M \to M\atop m \mapsto (\partial_{1} m, \partial_{2} m)\quad(m_1, m_2) \mapsto \partial_2(m_1) - \partial_1(m_2)$$On the other hand, we have $M\rightarrow \Omega^{1}_{X}\otimes_{R} M\rightarrow \Omega_{X}^{2}\otimes_{R}M$ to be $m\mapsto dx_{1}\otimes\partial_{1} m+dx_{2}\otimes\partial_{2}m$, $dx_{1}\otimes m_{1}+dx_{2}\otimes m_{2}\mapsto dx_{1}\wedge dx_{2}\otimes(\partial_{1} m_{2}-\partial_{2} m_{1})$
	
    The general case is similar. 
\end{proof}
	
\begin{example}
\quad
\begin{enumerate}
    \item Given $\alpha \in \mathbb{C}$, define $M_{\alpha} := \mathbb{C}[x, \frac{1}{x}] \cdot x^{\alpha}$, where $x^{\alpha}$ is a formal symbol such that $\partial_x x^{\alpha} = \alpha \frac{1}{x}\cdot x^{\alpha}$, then  $M_{\alpha}$ is a $\mathcal{D}_X$-module with: $$\partial_{x} \frac{1}{x^{k}} \cdot x^{\alpha} = (\alpha - k) \cdot \frac{1}{x^{k+1}} x^{\alpha} \quad \text{for } k \gg 0$$Since we obtain $\frac{1}{x^{k}}\cdot x^{\alpha}$ as a generator of $M_{\alpha}$ as $\mathcal{D}_{X}$-module, we have $M_{\alpha} \simeq \mathcal{D}_{\alpha} /\mathrm{Ann}_{\mathcal{D}_{\alpha}}( \frac{1}{x^{k}}\cdot x^{\alpha})\simeq\mathcal{D}_X / \mathcal{D}_X \cdot (x\partial_x - (\alpha - k))$. With $\mathrm{ch}(M)=V(\sigma(P))\subseteq T^{\ast}X$ in which $\sigma$ is the principal symbol, and $\sigma(P)=x\xi$, consider $T^{\ast}X\xrightarrow{\pi} X,\quad (x,\xi)\mapsto x$, we obtain $V(x)=\pi ^{-1}(0)=T^{\ast}_{\{0\}}X$ and $V(\xi)=T^{\ast}_{X}X$ zero section.  $\mathrm{ch}(M)=V(\sigma(x\partial_{x}))=T^{\ast}_{\{0\}}X\cup T_{X}^{\ast}X$ and  $$CC(M_{\alpha}) = T_{\{0\}}^*X + T_X^*X.$$Since $\mathrm{dim}ch(M)=1$ for $\mathrm{dim}T^{\ast}_{X}X=\mathrm{dim}X=1$, and $T^{\ast}_{\{0\}}X$ is the fiber of $T^{\ast}M$ at $0$, thus the cotangent space of $X$ at $0$, $\mathrm{dim}T^{\ast}_{\{0\}}X=\mathrm{dim}X=1$ as well. So  $M_{\alpha}$ is holonomic for all $\alpha \in \mathbb{C}$.
    \item Take $N = \mathbb{C}[\delta] \simeq \mathcal{D}_X / \mathcal{D}_X \cdot x \simeq \mathbb{C}[\partial_x]$, where $\delta$ is the Dirac delta function. Then $$CC(N) = T_{\{0\}}^*X$$also be holonomic. 

\end{enumerate}

\end{example}

Now, for $\alpha \notin \mathbb{Z}$, consider the de Rham complex of $M_{\alpha}$: $$M_{\alpha} \xrightarrow{\partial_1} M_{\alpha}$$

With $\text{Ker}(\partial_1) = \{f \cdot x^{\alpha} \mid \partial_1(f \cdot x^{\alpha}) = 0\}\simeq \{f \in \mathbb{C}[x, \frac{1}{x}] \mid xf' + \alpha f= 0\} = 0$, know that $x^{-\alpha}$ should be the solution of $xf'+\alpha f$, but $\alpha\not\in \mathbb{Z}$ makes $x^{-\alpha}\not\in \mathbb{C}\left[ x, \frac{1}{x} \right]$, thus has no solution in $\mathbb{C}[x, \frac{1}{x}]$. Yet, from ODE, we know that the differential equation $xf'+\alpha f=0$ has solutions at least locally, which suggests that the current definition of de Rham complex is not enough, and we need to extend scalars. 

\quad

Let $Y = \text{Spec } S$ be an affine variety of finite type over $\mathbb{C}$ with $\dim Y = n$. Define $Y \longrightarrow \mathbb{A}^{m}_{\mathbb{C}}$ by $$ \mathbb{C}[z_1, \ldots, z_m] \longrightarrow S\atop  z_i \mapsto \text{generators of } S$$And denote $I = I(Y) = \langle f_1, \ldots, f_k \rangle, \quad f_i \in \mathbb{C}[z_1, \ldots, z_m]$. 

In particular, $f_i$ are holomorphic functions on $\mathbb{C}^m$, then $f_1 = \cdots = f_k = 0$ defines an analytic subvariety $Y^{an}$ of the complex manifold $\mathbb{C}^m$. We have a continuous map: $$ \lambda: Y^{an} \longrightarrow Y $$For all$\ y \in Y^{\text{an}}$, by Implicit Function Theorem, $Y^{\text{an}}$ is smooth around $y$ if and only if the Jacobian matrix $J(f_1, \ldots, f_k)(y)$ at $y$ has rank $m-n$. 

\begin{definition} We define

\begin{enumerate}
    \item $Y$ is $\textbf{smooth}$ at a closed point $p$ if $$\dim_{\mathbb{C}} \mathfrak{m}_p/\mathfrak{m}_p^2 = n,$$ where $\mathfrak{m}_p$ is the maximal ideal of $p$.
    \item $Y$ is $\textbf{smooth}$ if $Y$ is smooth at every closed point.
\end{enumerate}
\end{definition}
\begin{lemma}
    $Y$ is smooth $\iff$ $Y^{\text{an}}$ is a complex manifold.
\end{lemma}
\begin{proof}
    Enough to prove for closed point $p$ of $Y$ we have: $$\text{rank}(J(f_1, \ldots, f_k)(p)) = m-n \quad (*)\iff \quad \dim_{\mathbb{C}} \mathfrak{m}_{p}/\mathfrak{m}_p^2 = n.$$Know that $p$ is also a closed point of $\mathbb{A}^{m}_{\mathbb{C}}$ with maximal ideal $\mathfrak{m} = (z_1 - \alpha_1, \cdots, z_m - \alpha_m)$. Define $$\mathbb{C}[z_1, \cdots, z_m] \to V = \mathbb{C}^m\atop f \quad \mapsto \quad \left( \frac{\partial f}{\partial z_1}(p), \cdots, \frac{\partial f}{\partial z_m}(p) \right)$$By HW, this induces an iso:$$
\mathfrak{\mathfrak{m}} /\mathfrak{m}^2 \xrightarrow[\varphi]{\simeq} V\simeq \Omega_{V,p}=\Omega_{V}\otimes R/\mathfrak{m}_{p}$$Since $I\subseteq \mathfrak{m}$ is generated by $f_{1},f_{2},\dots,f_{k}$, and for $h\in \mathbb{C}[z_{1},\dots,z_{m}]$, we can write $h=c+\sum(z_{i}-\alpha_{i})+\sum(z_{i}-\alpha_{i})(z_{j}-\alpha_{j})+\dots$, then we get $[hf_{i}]=c[f]$ in $\mathfrak{m} /\mathfrak{m}^{2}$ for $\mathfrak{m}^{2}$ makes only monomial with $\mathrm{deg}\leq 1$ matter and $f=\sum(z_{i}-\alpha_{i})+(\text{higher deg})$. So we have $\varphi(f_{1}),\dots,\varphi(f_{k})$ span $\varphi(I)=I+\mathfrak{m} /\mathfrak{m}^{2}$, and further, we obtain: $$
\dim(\varphi(I)) = \operatorname{rank} \left( J(f_1, \cdots, f_k)(p) \right).
$$
On the other hand, we have short exact sequence $$0\rightarrow\varphi(I)=I+\mathfrak{m}^{2} /\mathfrak{m}^{2}\rightarrow \mathfrak{m} /\mathfrak{m}^{2}\rightarrow \mathfrak{m}/\mathfrak{m}^{2}+I\rightarrow 0.$$
Know that $\mathfrak{m}_{p}=\mathfrak{m} /I$, $\mathfrak{m}_{p} /\mathfrak{m}_{p}^{2}=\frac{\mathfrak{m}/I}{\mathfrak{m}^{2}+I /I}\simeq \mathfrak{m} /\mathfrak{m}^{2}+I$,  with $\mathrm{dim}_{\mathbb{C}} \mathfrak{m} /\mathfrak{m}^{2}=m$, we get $\operatorname{rank} \left( J(f_1, \cdots, f_k)(p) \right)+\mathrm{dim}_{\mathbb{C}}\mathfrak{m}_{p} /\mathfrak{m}_{p}^{2}=m$, done. 
\end{proof}
\quad

Let $Z$ be a complex manifold (or variety) and $\mathcal{O}_{Z}$ the sheaf of all holomorphic functions on $Z$. 

We consider $Z=Y^{an}$, have $\lambda:Y^{an}\rightarrow Y$ induces a morphism of sheaves of rings $S \longrightarrow \mathcal{O}_{Y^{an}}$. If $M$ is an $S$-module, then $$
M^{\text{an}} \coloneqq M \otimes_S \mathcal{O}_{Y^{an}}
$$Then we finally define
$$
\operatorname{DR}(M) = K(M^{\text{an}}; \partial_1, \cdots, \partial_n)
$$
And $\mathrm{DR}(M)=[0 \rightarrow M^{an} \rightarrow M^{an} \otimes_R \Omega^1_X \rightarrow \cdots \rightarrow M^{an} \otimes_R \Omega^n_X \rightarrow0]=[0 \rightarrow M^{an} \rightarrow M^{an} \otimes_{\mathcal{O}} \Omega^1_{X^{an}} \rightarrow \cdots \rightarrow M^{an} \otimes_{\mathcal{O}} \Omega^n_{X^{an}}  \rightarrow0]$. 

Now, let's calculate $\mathrm{DR}(M_{\mathfrak{\alpha}})$. 

\begin{example}
    $X = \mathbb{A}^1_{\mathbb{C}}$,  $M_{\alpha} := \mathbb{C}[x, \frac{1}{x}] \cdot x^{\alpha}$.
    \begin{enumerate} 
    \item $\alpha \not\in \mathbb{Z}$, $\mathrm{DR}(M_{\alpha}) \cong [M_{\alpha} \xrightarrow{d_{\alpha}} M_{\alpha}]$. 
$\text{Ker}\partial_{x} \mid_{x \backslash \{0\}} \\ = \mathbb{C} \cdot x^{-\alpha} \cdot x^{d}$ since $\partial_{x}(x^{-\alpha}\cdot x^{\alpha})=\partial_{x}\cdot 1=0$. This implies $\text{Ker}\partial_{x} \mid_{x \backslash \{0\}}$ is the (non-trivial) local system $L_{\lambda}$ given by the multivalued function $x^{-\alpha}$, i.e. $T \cdot x^{-\alpha} = T \cdot e^{-\alpha \log x}= e^{-\alpha(\log x + 2\pi i)} = e^{-\alpha \cdot 2\pi i} e^{-\alpha \log x}= \lambda e^{-\alpha \log x}$, here $\lambda=e^{-\alpha2\pi i}$. 
$\partial_{x} \mid_{x \backslash \{0\}}$ is surjective, and then $\mathrm{DR}(M_{\alpha}) \mid_{x \backslash \{0\}} \simeq L_{\lambda}$, and  $H^{0}\mathrm{DR}(M_{\alpha}) \mid_{\{0\}} = H^{1}DR(M_{\alpha}) \mid_{\{0\}} = 0$. 
\item $\alpha \in \mathbb{Z} \Rightarrow \lambda = 1$, $\mathrm{DR}(M_{\alpha}) \mid_{x \backslash \{0\}} = \mathbb{C}_{x \backslash \{0\}}$, $H^{0}\mathrm{DR}(M_{\alpha}) \cong H^{1}\mathrm{DR}(M_{\alpha}) \cong \mathbb{C}_{\{0\}}$
\item $N=\mathbb{C}[\delta]\simeq \mathbb{C}[\partial_{x}]$, $\mathrm{DR}(N)\simeq \mathbb{C}_{\{0\}}[-1]$. 
    \end{enumerate}
\end{example}

\section{Basic version Riemman-Hilbert correspondence}

Let $Y$ be a complex manifold, $\mathcal{O}_Y$: the sheaf of holomorphic functions on $Y$, $\mathcal{L}$: a sheaf of $\mathbb{C}$-vector spaces. 

\begin{definition}
    $\mathcal{L}$ is called a $\textbf{local system}$ if $\forall y \in Y$, $\exists$ open neighborhood $U_y$ of $y$ s.t. $\mathcal{L}|_{U_y} \simeq\mathbb{C}^r_{U_y}$, where $\mathbb{C}^{r}_{U_y}$ is the constant sheaf on $U_y$.
\end{definition}

\begin{theorem}
    Let $X = \text{Spec } R$ be a smooth affine algebraic variety over $\mathbb{C}$ and $M$ an integrable connection on $X$, then $$M^{\text{an}} \simeq \mathcal{L} \otimes \mathcal{O}_{X^{\text{an}}},$$with $\mathcal{L}$ a local system.
\end{theorem}
\begin{proof}
   \quad

    Goal: $M$ is locally free, we want to show for small enough neighborhood, we can pick a basis $s_{1},\dots,s_{r}$ makes $\partial_{x_{i}}s_{j}=0$, thus
	We have known $M$ is a locally free $R$-module, $\forall x \in X$, take an open neighborhood $U \subset X$ s.t. $M|_U$ is free, then $M^{\text{an}}|_{U^{\text{an}}}$ is a free $\mathcal{O}_{U^{\text{an}}}$-module, generated by $\underline{\mathfrak{m}} = (m_1, \cdots, m_r)$.

	Denote $(x_1, \cdots, x_n)$ coordinates of $U^{\text{an}}$, then we have $\partial_{n} \underline{\mathfrak{m}} = X \cdot\underline{ \mathfrak{m}}$, where $X$ is a matrix of holomorphic functions on $U^{an}$. Expand $X$ w.r.t. $x_n$ in a small enough neighborhood of $x\in V\subseteq U^{an}$, we have $$X = \sum_{i\ge0} x_n^i \cdot X_i,$$where $X_i$ is a matrix of holomorphic functions in $(x_1, \cdots, x_{n-1})$. 

	Take $\Psi=\sum\limits_{i\ge 0} \frac{1}{i+1} x_{n}^{i+1}X_{i}$, so $\partial_{x}(\Psi)=X$, then $\Phi=\mathrm{\exp}(-\Psi)$, which is an invertible matrix. Set $\underline{S}=(s_{1},\dots,s_{r})=\Phi(\underline{\mathfrak{m}})$, so $\partial_{x}\underline{S}=\partial_{n}(\Phi)\cdot \underline{\mathfrak{m}}+\Phi \cdot(\partial_{n} \underline{\mathfrak{m}})=(\partial_{n}(\Phi)+\Phi \cdot X)\underline{\mathfrak{m}}=0$. 

	Since $\Phi$ is invertible matrix, we obtain $s_{1},\dots,s_{r}$ also a basis of $U^{an}|_{V}$, so we have $M^{an}|_{V}=\bigoplus\limits_{i=1}^{r}\mathcal{O}_{V}s_{i}$ s.t. $\partial_{n}s_{i}=0$. 

	Want to clarify the relation of $s_{i}$ and other $\partial_{x_{j}},j<n$, then Take $M_1 = \ker[M^{\text{an}}|_V \xrightarrow{\partial_n} M^{\text{an}}|_V]=\left\{ \partial_{n}\left( \sum f_{i}s_{i} \right)=0:f_{i}\in \mathcal{O}_{V} \right\}=\left\{ \sum\limits_{i=1}^{n}\partial_{n}(f_{i})s_{i}=0 \right\}.$ So we get $\sum_{i=1}^r \partial_n(f_i) s_i = 0$ for all $i=1,\dots,r$, i.e. $f_{i}$ is independent of $x_{n}$. 

	So we have $M_{1}=\bigoplus\limits_{i=1}^{r}\mathcal{O}_{V}^{n-1}s_{i}$, $\mathcal{O}_{V}^{n-1}$ are holomorphic functions that are independent of $x_{n}$, which suggests that we do not need the var $x_{n}$. And do the same thing for $M_{1}$,we get a new basis with $\partial_{n}s_{i}=\partial_{n-1}s_{i}=0$ for all $i=1,\dots,r$, and the change of basis is via an invertible matrix of homomorphic functions in $(x_{1},\dots,x_{n-1})$. 

	Continue the process, by induction, we finally obtain:  $$ M_n = \bigoplus_{i=1}^r \mathbb{C} \cdot s_i $$with $\partial_j(s_i) = 0 \quad \forall i, j$. And since every change of basis is by an invertible matrix, we know that for a small enough neighborhood $V \subseteq U^{an}$, we get $M^{an}|_{V}=\bigoplus\limits_{i=1}^{r}\mathcal{O}_{V}s_{i}$, with $\partial_{j}(s_{i})=0$ . 
\end{proof}

\begin{corollary}
    $\mathrm{DR}(M)|_{V}\stackrel{\text{quasi-isomorphism}}{\simeq}H^{0}\mathrm{DR}(M)|_{V}\simeq \mathbb{C}^{r}_{V}$. 
\end{corollary}
\begin{example}[Algebraic vs Analytic]
    Let $M = \mathbb{C}[x, \frac{1}{x}]$, $M' = \mathbb{C}[x, \frac{1}{x}] \cdot e^{\frac{1}{x}}$. 
    
    Know that $M$ and $M'$ are different as algebraic $D$-modules, but $M^{\text{an}}|_{\mathbb{C}^*} = \mathcal{O}_{\mathbb{C}^*}$, $M'^{an}|_{\mathbb{C}^{\ast}}=\mathcal{O}_{\mathbb{C}^{\ast}}\cdot e^{1/x}=\mathcal{O}_{\mathbb{C}^{\ast}}$, where $\mathbb{C}^{\ast}$ is the punctured complex plane. So $M^{\text{an}}|_{\mathbb{C}^*} \cong M'^{\text{an}}|_{\mathbb{C}^*}$ as analytic $\mathcal{D}$-modules and $DR(M^{\text{an}}) = DR(M'^{\text{an}}) \cong \mathbb{C}_{\mathbb{C}^*}^*$ Moreover, $M^{\text{an}} = \mathcal{O}_{\mathbb{C}}\cdot \frac{1}{x},M'^{an}=\mathcal{O}_{\mathbb{C}}\cdot e^{1/x}$ so $M^{an}$ and $M'^{an}$ are different as analytic $\mathcal{D}$-module. 
\end{example}

\section{Regular holonomic $\mathcal{D}$-module}

Let $X = \text{Spec } R$ be a smooth affine algebraic variety over $\mathbb{C}$ and $M$ a holonomic $\mathcal{D}_X$-mod. 
For any point $x \in X^{\text{an}}$, we consider the embedding:$$
i: \Delta \hookrightarrow X^{\text{an}}, \quad 0 \mapsto x.
$$where $\Delta$ is a small disk in the complex plane.
Then we have a morphism of sheaves of rings:$$R \to \mathcal{O}_{\Delta}, \quad f \mapsto f|_{\Delta}.$$We define $i^*M = M \otimes_R \mathcal{O}_{\Delta}$. 
\begin{definition}
    $M$ is called $\textbf{regular holonomic}$, if $i^{\ast}M$ is generated by $M_{\alpha}^{an}$ and $N^{an}|_{\Delta}$ in $\mathbf{K}_{0}(\mathcal{D}_{\Delta})$ for every small enough disk embedding $i:\Delta \hookrightarrow X^{an}$. 
\end{definition}

\begin{example}
    $M=\mathbb{C}\left[ x, \frac{1}{x} \right]\cdot e^{1/x}$ is not regular because $M$ is simple and $M\simeq \mathcal{D}_{X} /\mathcal{D}_{X}(x^{2}\partial_{x}+1)$. 
\end{example}


\section{Kashiwara's equivalence}
\begin{theorem}[Kashiwara's equivalence]
    Let $Y$ be a smooth affine alg. variety over $\mathbb{C}$  with coordinate $(x_{1},\dots,x_{n},y_{1},\dots,y_{\ell})$, and $X$ a subvariety defined by $(y_{1}=y_{2}=\cdots=y_{\ell}=0)$, so $X$ is smooth with coordinate $(x_{1},\dots,x_{n})$. Let $i: X \hookrightarrow Y$ be the closed embedding, $\mathrm{Mod}(\mathcal{D}_{X})$ the category of $\mathcal{D}_X$-mod and $\mathrm{Mod}^{X}(\mathcal{D}_{Y})$ the category of $\mathcal{D}_Y$-mod supported on $X$, then we have $$\mathrm{Mod}(\mathcal{D}_{X}) \xrightarrow{\simeq} \mathrm{Mod}^{X}(\mathcal{D}_{Y})\atop M\mapsto i_{+}M=M[\partial_{y_{1}},\dots,\partial_{y_{\ell}}].$$
\end{theorem}

\begin{proof}
    \quad

    By induction, it is enough to assume $$i: X \longrightarrow Y\atop (x_{1},\dots,x_{n}) \longrightarrow (x_{1},\cdots,x_{n},0)$$i.e. $X$ is defined by $\langle y = 0 \rangle$ in $Y$.

	Here we have $i_{+}: M\mapsto M[\partial_{y}]$, and we want to define a quasi-inverse of $i_{+}$, to define $i_{+}$, we shall study that how $\partial_{y}$ acts on $N$. 

	For $N \in \mathrm{Mod}^{X}(\mathcal{D}_{Y})$, consider the eigenvectors $N^{j} = \{ s \in N \mid y\partial_{y} \cdot s = j s \}$, want to show that $N=\bigoplus\limits_{i=1}^{\infty}N^{-i}$. 

	Since we have the relation $[\partial_{y}, y] = \partial_{y}y -  y\partial_{y} = 1$, we can check that $y N^{j} \subseteq N^{j+1}$, $\partial_{y} N^{j} \subseteq N^{j-1}$. And we have $y\partial_{y}: N^{j}\rightarrow N^{j}$ is isomorphism for all $j \neq 0$,  $\partial_{y}y = y\partial_{y} + 1: N^{j}\rightarrow N^{j}\quad\partial_{y}y\cdot s=(j+1)s$ is isomorphism for all $j \neq -1$. So for $j\neq -1$ and $0$, especially for $j<-1$, we have $\partial_{y}:N^{j+1}\rightarrow N^{j}$ and $y:N^{j}\rightarrow N^{j+1}$ are both isomorphisms. 

	Now we show $N=\bigoplus\limits_{i=1}^{\infty}N^{-i}$. 
	Since $N$ is supported on $X$, by Nullstellensatz, $\forall s\in N$ where $y^{k}\cdot s = 0$ for some $k\gg0$, i. e. $N = \bigcup\limits_{k\geq1} \ker(N\xrightarrow{y^{k}} N)$. 

	Since $N^{-i}$ are eigenspaces, $\bigoplus\limits_{i=1}^{\infty }N^{-i}\subseteq N$, suffices to prove that $N\subseteq \bigoplus\limits_{i=1}^{\infty }N^{-i}$. To show this, it's enough to prove that for all $k\ge 1$, we have $\ker(N\xrightarrow{y^{k}} N)\subseteq\bigoplus\limits_{i=1}^{k}N^{-i}$, and we do induction on $k$. 

	For $k=1$, $s\in \mathrm{ker(N\xrightarrow{y}N)}$, so we have $ys=0$, then $\partial_{y}ys=(y\partial_{y}+1)s=0,y\partial_{y}=-1\cdot s$, $\implies S\in N^{-1}$. 

	For $k>1$, now take $s\in \ker(N \xrightarrow{y^{k}}N)$, so we have $y^{k}s=y^{k-1} \cdot(ys)=0$, then $ys \in \ker(N \xrightarrow{y^{k-1}}N)$, which suggests $ys \in \bigoplus\limits_{i=1}^{k-1} N^{-i}$ by inductive hypothesis. 
	Since we have $\partial_{y}: N^{j+1}\rightarrow N^{j}$ , we obtain $\partial_{y}ys \in \bigoplus\limits_{i=2}^{k} N^{-i}$, and $\partial_{y}ys=y\partial_{y}s+s$. 

	On the other hand, we obtain  $\partial_{y}y^{k}s=y^{k}\partial_{y}s+ky^{k-1}s=y^{k-1}(y\partial_{y}s+ks)=0$, so  $y\partial_{y}s+ks\in \bigoplus\limits_{i=1}^{k-1} N^{-i}$ by inductive hypothesis. 

	Thus, $(k-1)s=(y\partial_{y}s+ks)-(y\partial_{y}s+s)\in \bigoplus_{j=1}^{k} N^{-i}$  implies that $s\in \bigoplus_{j=1}^{k} N^{-j}$. 

	With above, we now obtain $N=\bigoplus\limits_{i=1}^{\infty}N^{-i}$ with $N^{-1}=\mathrm{ker}(N\xrightarrow{y}N)$, so we have $N\simeq N^{-1}[\partial_{y}]$, and then we define $i^{\#}N = \ker(N \xrightarrow{y} N)$, $i^{\#}$ is the quasi-inverse of $i_{+}$. 

	In general, $i^{\#}N = H^{0}(K(N; y_{1}, \ldots, y_{k}))$.
\end{proof}

\begin{corollary}
    For $N \in \text{Mod}^{X}(\mathcal{D}_{Y})$ finitely-generated, we have $\text{Ch}(N) = T_{X}^{*}Y$ if and only if $i^{\#}N \in \text{Mod}(Q_{f})$ is an integrable connection.
\end{corollary}

\chapter{Auslander regularity and more on filtered rings$\&$modules}

\section{Auslander regular ring}
Let $A$ be a noetherian ring (not necessarily commutative), and $\text{Mod}_{f}(A)$, $\text{Mod}_{f}^{op}(A)$ the abelian category of left/right finitely generated $A$-module. For $M \in \text{Mod}_{f}(A)$, the $\textbf{projective dimension}$ of $M$ is $\mathrm{pd}(M) =$min of length of projective resolution. We say $A$ is $\textbf{of finite global dimension}$, if $\exists w \in \mathbb{Z}_{>0}$ s.t for all $M \in \text{Mod}_{f}(A)$, have $\mathrm{pd}(M) \le w$. We now assume that $A$ is of finite global dimension. Then $\text{RHom}(M,A) \in \mathbf{D}^{b}(\text{Mod}^{op}_{f}(A))$ the bounded derived category of $\mathrm{Mod}_{f}^{op}(A)$. 

\begin{definition}
\quad
    \begin{enumerate}
        \item The $\textbf{grade number}$ $j(M) := \min\left\{k \mid \text{Ext}^{k}(M,A) \neq 0\right\}$, we sometimes denote it as  $j_{A}(M)$ as well. In particular, define $j(M) = +\infty$ if $M = 0$. 
        \item For a noetherian and of finite global dimension ring $A$, we say $A$ is $\textbf{Auslander regular}$ (AR for short) if for all $M \in \text{Mod}_{f}(A)$ and $v \in \mathbb{Z}_{\ge 0}$, given any $N \subseteq \text{Ext}^{v}(M,A)$ submodule, have $j(N) \ge v$. 
    \end{enumerate}
\end{definition}

Let $S$ be a commutative $k$-algebra of finite type, we say $S$ is $\textbf{regular}$ if $\Omega_{S/k}$ is locally free. 

\begin{theorem}[Auslander]\label{theorem 7.2}
    Let $S$ be regular, $M \in \text{Mod}_{f}(S)$, then 
\begin{enumerate}
    \item $S$ is of finite global dimension; 
    \item $j(M) + \dim(\text{supp}(M)) = m$; 
    \item $\dim(\text{supp}(\text{Ext}^{i}(M,S))) \le m - i$ for all $i$; 
    \item $\dim(\text{supp}(\text{Ext}^{j(M)}(M,S))) = \dim(\text{supp}(M))$.
\end{enumerate}
\end{theorem}
\begin{proof}
    See Eisenbud for instance.
\end{proof}
\begin{corollary}
    Regular commutative $k$-algebra of finite type are Auslander regular.
\end{corollary}

Let $X = \text{Spec } R$ be a smooth-affine variety over $\mathbb{C}$, our next goal is to prove $\mathcal{D}_{X}$ is AR.

\section{Spectral sequences of filtered complex}

Let $F^{\bullet}$ be a bounded complex of $\mathbb{Z}$-modules:$$\cdots \to F^{i-1} \xrightarrow{d} F^i \xrightarrow{d} F^{i+1} \to \cdots$$with an increasing $\mathbb{Z}$-filtration: $F_j^i \subseteq F_{j+1}^i \subseteq F^i \quad \forall j,i \in \mathbb{Z}$ satisfying: 
\begin{enumerate}
    \item Exhaustive condition $\bigcup\limits_j F_j^i = F^i \quad \forall i$ ;
    \item Hausdorff condition $\bigcap_j F_j^i = 0 \quad \forall i$ 
    \item  Differential preserves filtration $d(F_j^i) \subseteq F_j^{i+1} \quad \forall i,j$ , which implies that $\dots\rightarrow F_{j}^{i-1}\xrightarrow{d} F_{j}^{i}\xrightarrow{d} F_{j}^{i+1}\rightarrow\dots$ is a complex. 
\end{enumerate}
Then, to every triple $(k, i, j)$ we set: 
\begin{itemize}
    \item $Z_j^i(k) = \{u \in F_j^i \mid du \in F^{i+1}_{j-k}\}.$
    \item $B^{i}_{j}(k)=F^{i}_{j}\cap d(F^{i-1}_{j+k-1})=d(Z^{i-1}_{j+k-1}(k-1))$
    \item $E_j^i(k) = \dfrac{Z_j^i(k) + F_{j-1}^i}{B_j^i(k) + F_{j-1}^i}$, and $E_{k}^{-j,i+j}:=E_{j}^{i}(k)$. And $E_{\bullet}^i(k) = \bigoplus\limits_{j} E_j^i(k)$.     
\end{itemize}

\begin{remark}
    Since we consider the increasing filtration $\dots\subseteq F_{j}\subseteq F_{j+1}\subseteq \dots \subseteq F$, while usually we consider a decreasing filtration for cochain complex, we can actually view $F_{j}$ as $F^{-j}$. That's the reason why we define $E_{k}^{-j,i+j}:=E_{j}^{i}(k)$.
\end{remark}
So $d(Z_j^i(k)) \subseteq Z_{j-k}^{i+1}(k)$, which induces $$ E_k^{-j,i+j} := E_j^i(k) \xrightarrow{d_k} E_k^{-j+k,i+j-k+1} = E_{j-k}^{i+1}(k) .$$Further, we have $$\dots\xrightarrow{d_{k}}E_{j+k}^{i-1}(k)\xrightarrow{d_{k}}E_{j}^{i}(k)\xrightarrow{d_{k}}E_{j-k}^{i+1}(k)\xrightarrow{d_{k}} \dots$$Consider $E^{i}_{\bullet}(k)=\bigoplus\limits_{j}E^{i}_{j}(k)$ and know that $d_{k}:E^{i}_{\bullet}(k)\rightarrow E^{i+1}_{\bullet}(k)$, this makes $E_\bullet(k)$ a complex of graded $\mathbb{Z}$-modules.  And $E_{j}^{i}(0) = F^{i}_{j} /F_{j-1}^{i}$, i.e. $E^{i}_{\bullet}(0)=\text{gr }_{\bullet}  F^{i}$ the associated graded complex of  $F^{i}$ for all $i$. 

$\{E_k^{-\ast,\bullet+\ast} = E_\ast^\bullet(k)\}_k$ is called the $\textbf{spectral sequence}$ of the filtered complex $F^{\bullet}$. 

\begin{lemma}\label{lemma 7.4}
    $E_\bullet^i(k+1) = H^i(E_{\bullet}^{\bullet}(k))$, more explicitly: $$E^{i}_{j}(k+1)\simeq \frac{\ker(E_{j}^{i}(k)\xrightarrow{d_{k}}E_{j-k}^{i+1}(k) )}{\mathrm{Im}(E_{j+k}^{i-1}(k)\xrightarrow{d_{k}}E_{j}^{i}(k))}.$$
\end{lemma}
\begin{proof}
    \quad

    First of all, consider $k=0$, we obtain that $E^{i}_{j}(0)= F^{i}_{j} /F^{i}_{j-1}$, so $\mathrm{ker}(E^{i}_{j }(0)\rightarrow E^{i+1}_{j}(0))=\ker(F^{i}_{j} /F^{i}_{j-1}\rightarrow F^{i+1}_{j} /F^{i+1}_{j-1})=\{u\in F^{i}_{j}:du\in F^{i+1}_{j-1}\}+F^{i}_{j-1} /F^{i}_{j-1}$ and $\mathrm{Im}(E^{i-1}_{j}(0)\rightarrow E^{i}_{j}(0))=\mathrm{Im}(F^{i-1}_{j} /F^{i-1}_{j-1}\rightarrow F^{i}_{j} /F^{i}_{j-1})=d(F^{i-1}_{j})+F^{i}_{j-1} /F^{i}_{j-1}$, then we get $E^{i}_{j}(1)\simeq \frac{\mathrm{ker}(E^{i}_{j }(0)\rightarrow E^{i+1}_{j}(0))}{\mathrm{Im}(E^{i-1}_{j}(0)\rightarrow E^{i}_{j}(0))}$. 

	For $k\ge 1$, we obtain that $F_{j-1}^{i}\subseteq F_{j+k-1}^{i}$$$E^{i}_{j}(k+1)=\frac{\{u\in F^{i}_{j}:du\in F^{i+1}_{j-k-1}\}+F^{i}_{j-1}}{F^{i}_{j}\cap d(F^{i-1}_{j+k})+F^{i}_{j-1}}.$$
    
    For $Z^{i}_{j}(k)+F^{i}_{j-1} /F^{i}_{j-1}\xrightarrow{d_{k}}Z_{j-k}^{i+1}(k)+F^{i+1}_{j-k-1} /F^{i+1}_{j-k-1}$, $\ker d_{k}:=\{u\in F^{i}_{j} /F_{j-1}^{i}: \exists x\in u\ s.t. dx\in F^{i+1}_{j-k-1}\}=Z^{i}_{j}(k+1)+F^{i}_{j-1} /F^{i}_{j-1}$. And $\mathrm{Im}(Z^{i-1}_{j+k}(k)+F^{i-1}_{j+k-1} /F^{i-1}_{j+k-1}\xrightarrow{d_{k}}Z^{i}_{j}(k)+F^{i}_{j-1}(k) /F^{i}_{j-1}(k))=d(Z^{i-1}_{j+k}(k))+F^{i}_{j-1}(k) /F^{i}_{j-1}(k)$, done. 
\end{proof}

\quad

The filtration of $F^{\bullet}$ induces a filtration of $H^i(F^{\bullet})$ by $$F_j H^i(F^\bullet) = \frac{Z_j^{i}(\infty) + dF^{i-1}}{dF^{i-1}}.$$where $Z_j^{i}(\infty) = \text{Ker}(d) \cap F_j^i$. 

In fact, we often view $Z^{i}_{j}(k)$ and $B^{i}_{j}(k)$ in $F^{i}_{j} /F^{i}_{j-1}$. Namely, we define $\eta^{i}_{j}:F^{i}_{j}\mapsto F^{i}_{j} /F^{i}_{j-1}$, and $\overline{Z}^{i}_{j}(k):=\eta^{i}_{j}(Z^{i}_{j}(k)),\overline{B}^{i}_{j}(k):=\eta^{i}_{j}(B^{i}_{j}(k))$, we obtain that $E^{i}_{j}(k)\simeq \frac{\overline{Z}^{i}_{j}(k)}{\overline{B}^{i}_{j}(k)}$. Moreover, we have a tower for all $i,j:$ $$
0 \subseteq \overline{B}_j^i(0) \subseteq \overline{B}_j^i(1) \subseteq \cdots \subseteq \overline{B}_j^i(r) \subseteq \cdots \subseteq \overline{B}_j^i(\infty)
\subseteq \overline{Z}_j^i(\infty) \subseteq \cdots \subseteq \overline{Z}_j^i(r) \subseteq \overline{Z}_j^i(r-1) \subseteq \cdots \subseteq \overline{Z}_j^i(0)
$$With $\overline{B}_j^i(\infty) =\varinjlim\limits_{r} \overline{B}_j^i(r) = \eta_j^i(F_j^i \cap d(F^{i-1}))$ (Using Exhaustive Condition for the second equality) and $\overline{Z}_j^i(\infty) = \varprojlim\limits_{ r} \overline{Z}_j^i(r) = \eta_j^i(Z_j^i(\infty))$ (Using Hausdorff Condition).  

\begin{proposition}
    $\mathrm{gr}_{\bullet}^{F}H^{i}(F)\simeq \frac{\overline{Z}^{i}_{j}(\infty)}{\overline{B}^{i}_{j}(\infty)}$. 
\end{proposition}

\quad

\begin{definition}
    We say the spectral sequence $\{E_{\bullet}^{\bullet}(k)\}_k$ $\textbf{converges to}$ $H^i(F^{\bullet})$ if for all $i,j$, there exists $k \gg 0$ s.t. $$E_j^i(k) \cong \text{gr}_j^F H^i(F^{\bullet}).$$
\end{definition}

\begin{definition}$\textbf{The regularity condition}:$
    if there exists a $\omega \in \mathbb{Z}_{\ge0}$ s.t.$$F_j^i \cap d(F^{i-1}_{}) \subset d(F_{j+\omega}^{i-1})$$for every pair $(i,j)$, then we say the filtered complex is $\omega$-regular.     
\end{definition}

\begin{lemma}\label{lemma 7.8}
    If filtered complex $F^{\bullet}$ is $\omega$-regular, then $$\mathrm{gr}^F_{\bullet} H^{i}(F^\bullet) \simeq E_{\bullet}^{i}(\omega+1), \quad \forall i$$In particular, the spectral sequence converges. 
\end{lemma}

\begin{proof}
    First of all, we show that $E^{i}_{j}(\omega+k)\simeq E^{i}_{j}(\omega+1)$ for all $k\ge 1$. 
	We have the long exact sequence:$$ \cdots \rightarrow E_{j+\omega+1}^{i-1}(\omega+1) \xrightarrow{d_{\omega+1}} E_j^i(\omega+1) \longrightarrow E_{j-\omega-1}^{i+1}(\omega+1) \rightarrow \cdots $$Suffices to show that every $d_{\omega+1}$ is zero map. 

	By Lemma \ref{lemma 7.4}, $E_j^i(\omega+1) = \left\{ \overline{u} \mid u \in E_j^i(\omega), d_\omega(\omega) = 0 \right\} = \left\{ \overline{u} \mid u \in F_j^i, du \in d\left(Z_{j-1}^i(\omega-1)\right) \right\}.$

    For any $\overline{\nu} \in d_{\omega+1}\left(E_{j+\omega+1}^{i-1}(\omega+1)\right)$, since the filtered complex is $\omega$-regular, we have $$ \nu \in F_j^i \cap d(F^{i-1}_{j+\omega+1}) \subseteq d(F_{j+\omega}^{i-1}), $$ then we get $E_j^i(\omega+1) = \frac{Z_j^i(\omega+1) + F_{j-1}^i}{B_j^i(\omega+1) + F_{j-1}^i} = \frac{Z_j^i(\omega+1) + F_{j-1}^i}{d\left(Z_{j+\omega}^{i-1}(\omega)\right) + F_{j-1}^i}$, where $d\left(Z_{j+\omega}^{i-1}(\omega)\right)=d\left( \left\{ u \in F_{j+\omega}^{i-1} \mid du \in F_j^i \right\} \right)$. 
	Thus we have $\overline{v}=0$ in $E^{i}_{j}(\omega+1)$, $E^{i}_{j}(\omega+k)\simeq E^{i}_{j}(\omega+1)$ for all $k\ge 1$, which implies that $\frac{\overline{Z}^{i+1}_{j}(\omega+k)}{\overline{B}^{i+1}_{j}(\omega+k)}\simeq\frac{\overline{Z}^{i+1}_{j}(\omega)}{\overline{B}^{i+1}_{j}(\omega)}$. 
	By Proposition, we obtain that $\frac{\overline{Z}^{i}_{j}(\infty)}{\overline{B}^{i}_{j}(\infty)}\simeq \mathrm{gr}_{\bullet}^{F}H^{i}(F)$, take $r\gg 0$, we consider the following map: 
    \begin{center}
        \begin{tikzcd}
            {\frac{\overline{Z}_j^i(r)}{\overline{B}_j^i(r)}} & {\frac{\overline{Z}_j^i(\infty)}{\overline{B}_j^i(r)}} & {\frac{\overline{Z}_j^i(\infty)}{\overline{B}_j^i(\infty)}} \\
             & {\frac{\overline{Z}_j^{i+1}(r+k)}{\overline{B}_j^{i+1}(r+k)}} 
             \arrow["\simeq"', draw=none, from=1-1, to=2-2]
             \arrow["\Phi"', hook', from=1-2, to=1-1] 
             \arrow["\Psi", two heads, from=1-2, to=1-3] 
             \arrow[from=1-2, to=2-2] 
             \end{tikzcd}
           
    \end{center}

	1. $\forall \overline{a} \in \frac{\overline{Z}_j^i(r)}{\overline{B}_j^i(r)}, \forall k > 0, \exists a \in \overline{Z}_j^i(r+k)$ s.t. $\overline{a} \in \frac{\overline{Z}_j^i(r+k)}{\overline{B}_j^i(r)}$ $\Rightarrow \overline{a} \in \frac{\overline{Z}_j^i(\infty)}{\overline{B}_j^i(r)} \Rightarrow \Phi$ is isomorphic. 
	
    2. $\forall b \in \frac{\overline{Z}_j^i(\infty)}{\overline{B}_j^i(r)}$ s.t. $\psi(\overline{b}) = 0$ $\Rightarrow0= \overline{\overline{b} }\in \frac{\overline{Z}_j^i(\infty)}{\overline{B}_j^i(r+k_0)}$ for some $k_0 \gg 0\Rightarrow \overline{\overline{b}} \in \frac{\overline{Z}_j^i(r+k_0)}{\overline{B}_j^i(r+k_0)}$, but $\overline{b} \in \frac{\overline{Z}_j^i(\infty)}{\overline{B}_j^i(r+k_0)} \Rightarrow \overline{b} = 0\Rightarrow \Psi$ is isomorphic. 

	Therefore, $E_j^i(k)\simeq\frac{\overline{Z}^{i}_{j}(\infty)}{\overline{B}^{i}_{j}(\infty)} \simeq \mathrm{gr}_j H^i(F) $ for all $j$. 
\end{proof}

\begin{theorem}\label{tehorem 7.9}
    If $F^\bullet$ is $\omega$-regular, then $\mathrm{gr}^F_{\bullet} H^{i}(F^{\bullet})$ is a subquotient of $H^{i} \mathrm{gr}^{F}_{\bullet }(F^{\bullet})$. 
\end{theorem}
\begin{proof}
    By Lemma \ref{lemma 7.4} + Lemma \ref{lemma 7.8}. 
\end{proof}

\section{More on filtered rings }

Let $A$ be a ring and $\Gamma_{\bullet} = \Gamma_{\bullet}A$ a $\mathbb{Z}$-filtration on $A$, recall that $\Gamma_{\bullet}A$ is called a $\textbf{noetherian } \mathbb{Z}\textbf{-filtration}$ if additionally $\mathrm{gr}^{\Gamma}_\bullet A$ is also noetherian.

Let $M$ be a finitely-generated $A$-module, and $\Omega_{\bullet}M$ a good filtration on $M$ s.t. $\mathrm{gr}^{\Omega} M \neq 0$. 
Since $\mathrm{gr}^{\Gamma} A$ is noetherian, $R^{\Omega} M$ has a graded free resolution: $$\dots\rightarrow P_{\bullet}^{-1}\rightarrow P^{0}_{\bullet}\rightarrow R^{\Omega} _{\bullet}M\rightarrow0$$s.t. each $P^{i}_{\bullet}$ is a free graded $R_{\bullet}^{\Gamma}A$-module of finite rank. 
And $\mathrm{RHom}_{R_{\bullet}A}(R_{\bullet}M, R_{\bullet}A)$ is the complex $$\mathrm{Hom}_{R_{\bullet}A}(P^{\bullet}_{\bullet}, R_{\bullet}A).$$
Tensoring $\mathbb{C}[u]/(u-1)$, we get a free resolution:
$$ \cdots \to F^{-1} \to F^0 \to M \to 0$$
with $F^i = P^i_{\bullet} \otimes \mathbb{C}[u]/(u-1)$.
Define the filtration on $F^i_{\bullet}$ by:$$F_j^i = \text{Im}(P_j^i \longrightarrow F^i)$$This filtered complex:$$\cdots \longrightarrow F_{\bullet}^{-1} \longrightarrow F_{\bullet}^0 \longrightarrow 0$$is called a $\textbf{filtered free resolution}$ of $M$.

Since $\mathrm{Hom}(\underline{},A)$ functor is contravariant, we define a filtered complex $$ \check{F}^{\bullet} = \operatorname{Hom}_{A}(F^{\bullet}, A) \quad \text{s.t.} \check{F}^{i} = \operatorname{Hom}_{A}(F^{-i}, A) $$We consider $\mathrm{Hom}_{R_{\bullet}A}(P^{\bullet}_{\bullet},R_{\bullet}A)\otimes \mathbb{C}[u] /(u-1)$, since $P^\bullet_{\bullet}$ is a free $R_{\bullet}A$-module, we then have $\mathrm{Hom}_{R_{\bullet}A}(P^\bullet_{\bullet},R_{\bullet}A)\otimes \mathbb{C}[u]/(u-1)\simeq \mathrm{Hom}_{A}(F^{-1},A)$ and filtration  $\check{F}^{i}_{j} = \operatorname{Im}\left( \operatorname{Hom}_{R_{\bullet}A}(P^{-i}_{j} , R_{\bullet}A)\rightarrow \check{F}^{i} \right)$. 

\begin{lemma}
    $\check{F}^{\bullet}$ is  $\omega$-regular for some $\omega \gg 0.$
\end{lemma}
\begin{proof}
    Take $N = d(\check{F}^{k}) \subseteq \check{F}^{k+1}$, then $N_j = d(\check{F}_{j}^{k})$ gives a good filtration on $N$.
	On the other hand, $N \cap \check{F}_{j}^{k+1}$ gives another good filtration. By Corollary 3.4, there exists a uniform $\omega$ s.t. $$ \check{F}_{j}^{k+1} \cap d(\check{F}^{k}) = N \cap \check{F}_{j}^{k+1} \subseteq N_{j+\omega} = d(\check{F}_{j+\omega}^{k}), \quad \forall j. $$
\end{proof}

\begin{corollary}\label{corollary 7.11}
    $\mathrm{gr}_{\bullet}\mathrm{Ext}_A^j (M, A)$ is a subquotient of $\mathrm{Ext}_{\mathrm{gr}_{\bullet} A}^j (\mathrm{gr} _{\bullet}M, \mathrm{gr}_{\bullet} A)$ for all $j$. 
\end{corollary}
\begin{proof}
    This is just a special case of Theorem $\ref{tehorem 7.9}$. 
\end{proof}

\section{Duality for $\mathcal{D}_{X}$-mod}

Let $X = \text{Spec } R$ a smooth affine variety with algebraic coordinates $(x_1, \ldots, x_n)$ and $M$ a nonzero f.g. $\mathcal{D}_X$-module. 

\begin{theorem}\label{theorem 7.12}
    We have:
    \begin{enumerate}
        \item $\dim \operatorname{ch}(\operatorname{Ext}_{\mathcal{D}_X}^i(M, \mathcal{D}_X)) \leq 2n - i$; 
        \item $\operatorname{Ext}_{\mathcal{D}_X}^i(M, \mathcal{D}_X) = 0$ for $i < 2n - \dim(\operatorname{ch}(M))$ or $i > n$; 
        \item $j(M) + \dim \operatorname{ch}(M) = 2n$. 
    \end{enumerate}
\end{theorem}
\begin{proof}
    1) and 2) holds for Theorem \ref{theorem 7.2} and Corollary \ref{corollary 7.11}. 

    $3)$ Know that $j(\mathrm{gr}_{\bullet}M)+\mathrm{dim}(Ch(M))=\mathrm{dim}R[\xi_{1},\dots,\xi_{n}]=2n$, suffices to show that $j(M)=j(\mathrm{gr}_{\bullet}M)$. Since $\mathrm{gr}_{\bullet}\mathrm{Ext}_{\mathcal{D}_{X}}^{j}(M,\mathcal{D}_{X})$ is a subquotient of $\mathrm{Ext}^{i}_{\mathrm{gr_{\bullet}\mathcal{D}_{X}}}(\mathrm{gr}_{\bullet}M,\mathrm{gr}_{\bullet}\mathcal{D}_{X})$, we obtain that $j(M)\ge j(\mathrm{gr}_{\bullet}M)=2n-\mathrm{dim}(Ch(M))$. 

	Assume $j(\mathrm{gr}_{\bullet}M)=j_{0}$, which is to say that $\mathrm{Ext}^{j_{0}}_{\mathrm{gr_{\bullet}A}}(\mathrm{gr}_{\bullet}M,\mathrm{gr}_{\bullet}\mathcal{D}_{X})\neq 0$, and $\mathrm{Ext}^{j}_{\mathrm{gr}_{\bullet}\mathcal{D}_{X}}(\mathrm{gr}_{\bullet}M,\mathrm{gr}_{\bullet}\mathcal{D}_{X})$ is the $E^{1}$-pages of spectral sequence that converges to $\mathrm{Ext}_{\mathcal{D}_{X}}^{j}(M,\mathcal{D}_{X})$. Also, $E^{i}_{\bullet}(0)=\mathrm{gr}_{\bullet}\check{F}^{i}$. 

	We consider $E^{j}_{\bullet}(1)\simeq \mathrm{Ext}^{j}_{\mathrm{gr}_{\bullet}\mathcal{D}_{X}}(\mathrm{gr}_{\bullet}M,\mathrm{gr}_{\bullet}\mathcal{D}_{X})$, and then for $j<j_{0}$, we then have $E^{j}_{\bullet}(1)=0$. Thus $E^{j}_{\bullet}(k+1)\simeq \mathrm{ker}(E^{i}(k)\rightarrow E^{i+1}(k))$ for $k\ge 1$. 

	Since we take the order filtration of $\mathrm{gr}_{\bullet}\mathcal{D}_{X}$, which is a $\mathbb{Z}_{\ge 0}$-filtration and a good filtration on $M$ up to $\mathrm{gr}_{\bullet}\mathcal{D}_{X}$, we have $\Omega_{p}M=\sum\limits_{\text{finite}}\Gamma_{p-k_{i}}\mathcal{D}_{X}\cdot m_{i}$ for some $m_{i}\in \Omega_{k_{i}}M$. We will get $\Omega_{p}M=0$ for $p<\min\{k_{i}\}$, i.e. $\Omega_{p}M=0$ for $p\ll0$. Same for the filtration induced by a free resolution of $\Omega_{p}M$, it must be bounded below. Take $p_{0}$ the minimal integer s.t. $\check{F}^{j_{0}}_{p_{0}}\neq 0$. We know that $E^{j_{0}}_{p_{0}}(0)=\check{F}^{j_{0}}_{p_{0}}\neq 0$ and for $k\ge 1$, $d_{k}:E^{j_{0}}_{p_{0}}(k)\rightarrow E^{j_{0}+1}_{p_{0}-k}(k)$ must be zero map. So $E^{j_{0}}_{p_{0}}(\infty)=E^{j_{0}}_{p_{0}}(0)\neq 0$, we which suggests that $\mathrm{gr}_{p_{0}}\mathrm{Ext}_{\mathcal{D}_{X}}^{j_{0}}(M,\mathcal{D}_{X})\neq 0$, and thus $\mathrm{Ext}_{\mathcal{D}_{X}}^{j_{0}}(M,\mathcal{D}_{X})\neq 0$, we then have $j(M)\leq j(\mathrm{gr}_{\bullet}M)$, done.  
\end{proof}

\begin{proposition}
    The global dimension of $D_X$ is also finite. Furthermore, the global dimension of $\mathcal{D}_X$ $=n$.
\end{proposition}
\begin{proof}
    We have $\mathrm{dim}X=n$ with coordinates $(x_{1},\dots,x_{n})$, since $X$ is smooth affine, we have $\Omega_{X /k}$ is free thus locally free. Then $X$ is regular with $\mathrm{gl\ dim}X=\mathrm{dim}X=n$. 

	Pick $M$ a finitely-generated $\mathcal{D}_{X}$-module, with $M$ is finitely-generated over $R$, or just take $M=\mathcal{O}_{X}$, know that $M$ is an integrable connection, thus $M$ is holonomic with $\mathrm{dim}(\mathrm{ch}(M))=n$, which indicates that $j(M)=2n-n=n$. Thus the $\mathrm{pd}(\mathcal{O}_{X})\ge n$, and $\mathrm{gl\ dim}\mathcal{D}_{X}\ge n$. 

	Claim that for $\mathcal{D}_{X}$, we obtain that $\mathrm{pd}(M)=j(M)$, and then since $j(M)=2n-\mathrm{dim}(\mathrm{ch}(M))\leq n$ , we have $\mathrm{gl\ dim}\mathcal{D}_{X}\leq n$. 

	To prove the claim, easy to check that $j(M)\leq \mathrm{pd}(M)$, since $j(M)\leq n$ is finite.  

	For $\mathrm{pd}(M)=0$, thus $j(M)=0$.  
	
    Assume for $\mathrm{pd}_{R}(M)<k$, we have $j(M)=\mathrm{pd}(M)$, now consider $\mathrm{pd}(M)=k$, we know that there exists a resolution of $M$ s.t. $$0\rightarrow P\rightarrow F_{k-1}\rightarrow \dots\rightarrow F_{0}\rightarrow M\rightarrow0$$ with $F_{i}$ are all free and $P$ is projective. we consider the exact sequence $$0\rightarrow K\rightarrow F_{0}\twoheadrightarrow M\rightarrow 0 ,$$where we take $K$ to be the kernel, then $\mathrm{pd}K=\mathrm{pd}M-1<k$. And we know that $\mathrm{Ext}^{i}(K,\mathcal{D}_{X})=\mathrm{Ext}^{i+1}(M,\mathcal{D}_{X})$, since $\mathrm{Ext}^{i}(K,\mathcal{D}_{X})=0$ for $i\neq \mathrm{pd}(K)=j(K)=k-1$, then $\mathrm{Ext}^{i}(M,\mathcal{D}_{X})=0$ for all $i\neq k$, thus $j(M)=k$, done. 
\end{proof}

\begin{theorem}
    $\mathcal{D}_{X}$ is Auslander regular. 
\end{theorem}
\begin{proof}
    If $N \subset \operatorname{Ext}^{i}(M, \mathcal{D}_X)$ is a nonzero submodule, then $\operatorname{ch}(N) \subset \operatorname{ch}(\operatorname{Ext}^{i}(M, \mathcal{D}_X))$. So $\dim(\operatorname{ch}(N)) \leq \mathrm{dim}(\operatorname{ch}(\operatorname{Ext}^{i}(M, \mathcal{D}_X)))\leq 2n - i$ by Theorem \ref{theorem 7.12}.1. Then $j(N) \geq i$ by Theorem \ref{theorem 7.12}.3. 
\end{proof}

\begin{corollary}
    \quad
    \begin{enumerate}
        \item $M$ is holonomic iff $\operatorname{Ext}^i(M, \mathcal{D}_X) = 0$, $i \neq n$; 
        \item If $M$ is holonomic then $\operatorname{Ext}^n(M, \mathcal{D}_X)$ is also holonomic.
    \end{enumerate}
\end{corollary}



\chapter{Algebraic normal deformation and V-filtrations}

\section{$V$-filtration and normal deformation}
Let $X = \text{Spec } R$ be a smooth affine variety over $\mathbb{C}$ with algebraic coordinates $(t, x_2, \ldots, x_n)$. 

Take $I^j =\begin{cases}\left< t \right>^{j},& j > 0\\R,&j\leq0\end{cases}$, we define a filtration of $\mathcal{D}_{X}$: $$V_{k} \mathcal{D}_X = \{P \in \mathcal{D}_X \mid P I^{j} \subseteq I^{j-k}\text{ for all }j\in \mathbb{Z}\}.$$Then we obtain that $V_0 \mathcal{D}_X = R \langle t\partial_t, \partial_{x_2}, \ldots, \partial_{x_n} \rangle$ and for all $i>0$: $$V_{-i} \mathcal{D}_{X} = t^{i} V_{0 }\mathcal{D}_{X }= V_{0} \mathcal{D}_{X} t^{i}\atop V_i \mathcal{D}_X = \partial_t V_{i-1} \mathcal{D}_X + V_{i-1} \mathcal{D}_X.$$Next, we will prove $V_{\bullet} \mathcal{D}_X$ is a noetherian $\mathbb{Z}$-filtration s.t. $\mathrm{gr}^V \mathcal{D}_X$ is Auslander regular. 

We obtain that $V_{j} \mathcal{O}_X = I^{-j} \in \mathcal{O}_X = R$ for $j \in \mathbb{Z}$ defines a $\mathbb{Z}$-filtration on $\mathcal{O}_X$.
We have 
$$\mathcal{R}_{\bullet}(\mathcal{O}_X) = \mathcal{R}^V_{\bullet}(\mathcal{O}_X) = \bigoplus_{j \in \mathbb{Z}} I^{-j} \otimes u^j \in R[u, u^{-1}]$$
and $\mathbb{C}[u] \hookrightarrow R_{\bullet}(\mathcal{O}_X)=R[u]\bigoplus\limits_{j\ge 1}I^{j}\otimes u^{-j}$ induces a family $\widetilde{X}= \operatorname{Spec} R_{\bullet}(\mathcal{O}_X)\xrightarrow{\varphi}\mathbb{A}^{1}_{\mathbb{C}}$ s.t. 
\begin{enumerate}
    \item $\varphi$ is smooth affine; 
    \item  $\varphi^{-1}(a) \simeq X, \quad \forall \text{ nonzero } a \in \mathbb{A}^1_{\mathbb{C}}$, since $\varphi ^{-1}(a)=X\times_{\mathbb{A}^{1}_{C}}\mathrm{Spec}(\mathbb{C}[u] /(u-a))=\mathrm{Spec}(R\otimes_{\mathbb{C}[u]}\mathbb{C}[u] /(u-a))=\mathcal{R}_{\bullet}(\mathcal{O}_{X}) /(u-a)\simeq R$. 
    \item $\varphi^{-1}(0) = \operatorname{Spec} \bigoplus\limits_{j \in \mathbb{Z}_{\geq 0}} I^j/I^{j+1} = T_Y X$ the normal bundle of $Y$ in $X$, where $Y = \left<  t = 0\right> \subseteq X$. 

\end{enumerate}

We say $\varphi$ is the $\textbf{normal deformation}$ along $Y$. 

Recall the following definition: 

\begin{definition}
    Let $A$, $B$ be two commutative rings with units with $\phi: A \to B$: a ring homomorphism. Then $\Omega_{\phi} = \Omega_{B/A}$ is the $B$-module generated by $df$, $f \in B$ s.t.
    \begin{enumerate}
        \item  $d(fg) = fdg + gdf$
        \item $d(af + bg) = adf + bdg$, $\quad \forall f,g \in B$, $a,b \in A$
    \end{enumerate}
 called the $\textbf{module of Kähler differentials}$ of $B$ over $A$. 

 Furthermore, we have the following definition: 
 \begin{enumerate}
    \item  $\phi$ is called $\textbf{smooth}$ if $\Omega_{\phi}$ is a free $B$-module of finite rank.
    \item If $\phi$ is smooth, we define $$\mathcal{T}_{\phi} := \mathrm{Hom}_B(\Omega_{\phi}, B),$$which is called the $\textbf{module of derivatives}$ of $B$ over $A$.
    \item $\mathcal{D}_{\phi} \subseteq \mathrm{End}_A(B)$ is the $A$-subalgebra generated by $\mathcal{T}_{\phi}$ and $B$, the $\textbf{ring of relative differential operators}$ of $B$ over $A$. 
 \end{enumerate}

\end{definition}

\begin{remark}
    if $A'$ is another $A$-algebra, then $\Omega_{B'/A'} \simeq \Omega_{B/A} \otimes_B B',$ where $B' = B \otimes_A A'$.
\end{remark}

\begin{lemma}
    
    Consider the ring homomorphism $\mathbb{C}[u]\hookrightarrow \mathcal{R}_\bullet(\mathcal{O}_X)$, we obtain that:

    \begin{enumerate}
        \item $\Omega_\varphi = \bigoplus_{j \in \mathbb{Z}} I^{-j} (dt \otimes u^{j-1}  \bigoplus\limits_{i=2}^{n} dx_{i} \otimes u^{j}) = \mathcal{R}_{\bullet}(\mathcal{O}_X) dt \otimes u^{-1}  \bigoplus\limits_{i=2}^{n} \mathcal{R}_{\bullet}(\mathcal{O}_X) dx_{i}.$
        \item $\mathcal{T}_\varphi = \bigoplus\limits_{j \in \mathbb{Z}} I^{-j} (\partial_t \otimes u^{j+1} \bigoplus_{i=2}^n \partial_i \cdot u^j)$
    \end{enumerate}
\end{lemma}
\begin{proof}
     Since $d(t^k \cdot u^{-k}) = t^{k-1} dt \cdot u^{-k}$, 1 follows. And taking graded dual, 2 follows.
\end{proof}

\begin{corollary}
    \quad
    \begin{enumerate}
        \item $\mathcal{R}_{\bullet}^V \mathcal{D}_X = \mathcal{D}_\varphi$
        \item  $\mathrm{gr}_{\bullet}^V \mathcal{D}_X \cong \mathcal{D}_{T_Y X}$
    \end{enumerate}
    Then $\mathcal{R}^V_{\bullet} \mathcal{D}_X$ is noetherian and  $\mathrm{gr}^{V}_{\bullet} \mathcal{D}_X$ is Auslander regular. Meanwhile, we obtain that $\bigcap\limits_{i \in \mathbb{Z}} V_i \mathcal{D}_X = 0$.  

\end{corollary}

\section{Bjork’s approach to existence of b-functions}

\begin{theorem}[Björk]\label{Theorem 8.4}
    Let $M$ be a holonomic $\mathcal{D}_X$-module with a good filtration $\Omega_{\bullet} M$ over $V_{\bullet} \mathcal{D}_X$. There exists $b(s) \in \mathbb{C}[s]$ s.t. $$b(t\partial_t - k) \text{ kills } \frac{\Omega_k M}{\Omega_{k-1} M}, \quad \forall k \in \mathbb{Z}$$The monic minimal degree $b(s)$ is called the $\textbf{b-function}$ of $M$ along $Y$.
\end{theorem}
\begin{lemma}\label{lemma 8.5}
    Let $Y$ be a smooth variety over $\mathbb{C}$, $N$ a holonomic $\mathcal{D}_Y$-module, and $\theta\in End_{\mathcal{D}_Y}(N)$; then there exists $b(s)\in\mathbb{C}[s]$ s.t. $b(\theta)=0$. 
\end{lemma}
\begin{proof}
    
    Since $N$ is holomonic, we obtain that$\mathrm{ch}(N) = \bigcup\limits_{i} T^*_{Y_i} Y$, with $Y_{i}$ locally closed. Take a $Y_i$ of maximal dimension and take open subset $U \subset Y$ s.t. $Y_{i}^{\circ} \hookrightarrow U$ a closed embedding s.t. $\overline{Y_i^{\circ}} = Y_{i}$ and $U \cap \{\bigcup\limits_{j\neq i}Y_j\} = \varnothing$, then $\mathrm{ch}(N|_U) = T^*_{Y_i^\circ} Y\Rightarrow N|_U = \mathcal{O}_{Y_i^\circ} \otimes L[\partial_1, \ldots, \partial_k]$ with $L$ a local system on $Y_i^\circ$. We have $\mathrm{End}(N|_U) = \mathrm{End}_{\mathbb{C}}(L)$, therefore, $\exists b_i(s) \in \mathbb{C}[s]  \text{ s.t. } b_i(\theta|_{U}) = 0.$

	Now, take $N_1 = b_1(\theta)N$, then $\mathrm{ch}(N_1) = \bigcup\limits_{j \neq i} T^*_{Y_j} X$, We are done by induction. 
\end{proof}
\begin{theorem}\label{Theorem 8.6}
    Let $M$ be a finitely generated $\mathcal{D}_{X}$-module, and $\Omega_{\bullet}M$ a good filtration over $V_{\bullet}\mathcal{D}_{X}$, then $j(\mathrm{gr}_{\bullet}^\Omega M)\ge j(M)$.  
\end{theorem}
\begin{proof}
    If $\mathrm{gr}^\Omega_{\bullet}M=0$, nothing to prove, so we assume $\mathrm{gr}^\Omega_{\bullet}M\neq 0$ and $\ell=j(\mathrm{gr}_{\bullet}^{\Omega}M)<j(M)$. 
	Then the $E_1$-page of the spectral sequence that converges to $\mathrm{Ext}_{\mathcal{D}_{X}}^{j}(M,\mathcal{D}_{X})$ satisfies:$$0 \rightarrow \operatorname{Ext}^{\ell}(\operatorname{gr}_{\bullet}^{\Omega} M, \operatorname{gr}^{V}_{\bullet}\mathcal{D}_X) 
\xrightarrow{d_1} 
\operatorname{Ext}^{\ell+1}(\operatorname{gr}_{\bullet}^{\Omega} M, \operatorname{gr}^{V}_{\bullet}\mathcal{D}_X)$$So we obtain that $E^{\ell}_{\bullet}(2)=\mathrm{ker}d_{1}\subseteq\operatorname{Ext}^{\ell}(\operatorname{gr}_{\bullet}^{\Omega} M, \operatorname{gr}^{V}_{\bullet}\mathcal{D}_X)$ since no images map to $E^{\ell}_{\bullet}(1)$, then we have the following exact sequence with the quotient $N=E^{\ell}_{\bullet}(1) /\ker d_{1}\simeq d_{1}(E_{\bullet}^{\ell}(1))\subseteq E^{\ell+1}_{\bullet}(1)$: $$0 \rightarrow E_{\bullet}^{\ell}(2) \rightarrow \operatorname{Ext}^{\ell}(\operatorname{gr}_{\bullet}^{\Omega} M, \operatorname{gr}^{V}_{\bullet}\mathcal{D}_X) 
\rightarrow N \rightarrow 0$$With $\mathrm{gr}_{\bullet}^V\mathcal{D}_{X}$ Auslander regular, we obtain that for $N\subseteq \mathrm{Ext}^{\ell+1}(\operatorname{gr}_{\bullet}^{\Omega} M, \operatorname{gr}^{V}_{\bullet}\mathcal{D}_X)$, $j(N)\ge \ell+1$. 
	Since $j(M)>\ell$, we have $E_{\bullet}^{\ell}(r) = \operatorname{gr}_{\bullet} \operatorname{Ext}^{\ell}(M, \mathcal{D}_X) = 0$ for some $r\gg 0$, and $E^{\ell}_{\bullet}(r)=\mathrm{ker}d_{r}.$

	Also $\mathrm{gr}_{\bullet}^V\mathcal{D}_{X}$ is a noetherian ring, we then have $j(\operatorname{Ext}^{\ell}(\operatorname{gr}_{\bullet}^{\Omega} M, \operatorname{gr}^{V}_{\bullet}\mathcal{D}_X) )\ge\ell+1$, a contradiction. 

\end{proof}


\begin{proof}[proof of (\ref{Theorem 8.4})]
   \quad
   
   By Theorem \ref{Theorem 8.6}, $j(\mathrm{gr}_{\bullet}^\Omega M)\ge j(M)=n$. If $\mathrm{gr}_{\bullet}^\Omega M=0$, take $b(s)=0$. If not, then $\mathrm{dim\ ch}(\mathrm{gr}^\Omega_{\bullet}M)\ge n$. 

   Since $\mathrm{gr}_{\bullet}^V\mathcal{D}_{X}$ is a Auslander regular ring, we obtain that $j(\mathrm{gr}_{\bullet}^\Omega M)+\mathrm{dim}(\mathrm{ch}(\mathrm{gr}_{\bullet}^\Omega M))=2n=2\mathrm{dim}T_{Y}X$, then $\mathrm{gr}^\Omega_{\bullet} M$ is holonomic over $T_Y X$. 

   Now, take $\theta=\bigoplus \limits_{i}(t\partial_{t}-i),$ then $\theta\in\mathrm{End}_{\mathcal{D}_{T_{Y}X}}(\mathrm{gr}_{\bullet}^\Omega M)$. By Lemma \ref{lemma 8.5}, there exists $b(s)\in\mathbb{C}[s]$ s.t. $b(\theta)=0$, In particular, $b(t\partial_{t}-k)$ kills $\frac{\Omega_{k}M}{\Omega_{k-1}M}$ for all $k$. 
\end{proof}

\begin{remark}
    It is poosible that $\mathrm{gr}_{\bullet}^\Omega M=0$. In fact, take $M=\mathbb{C}[t, \frac{1}{t}]e^{\frac{1}{t}}=\mathcal{D}_{X}\cdot e^{\frac{1}{t}}$, take $\Omega_{k}M=V_{k}\mathcal{D}_{X}\cdot e^{\frac{1}{t}}$, which is a good filtration. Since $t\partial_{t}e^{\frac{1}{t}}= -\frac{1}{t} e^{\frac{1}{t}}$, we have $\Omega_{k}M=M$ and $b(s)=1$. 
\end{remark}

\begin{theorem}[Kashiwara-Malgrange]\label{Theorem 8.7}
    Let $M$ be a holonomic $\mathcal{D}_{X}$-module, there exists unique good filtration $V_{\bullet}M$ and $b(s)\in\mathbb{C}[s]$ s.t. 
    \begin{enumerate}
     \item $\mathrm{Re}(\text{roots of }b(s))\subseteq(-1,0]$;
     \item $b(t\partial_{t}-k)$ kills $\frac{V_{k}M}{V_{k-1}M}$. 
    \end{enumerate}
    And the filtration is called $V$-filtration. 
\end{theorem}
\begin{proof}
    This is Kashiwara's trick: 
    Set $\Gamma_{n}=\Omega_{n+k}$, we then have $b(t\partial_{t}-n-k)$ kill $\Gamma_{n} /\Gamma_{n-1}$ for all $n$, thus taking $k\gg 0$, we can assume $\mathrm{Re}(\text{roots of }b(s))>0$. 
    
    Let $b(s)=(s-\alpha)^\ell b_{1}(s)$ s.t. $(s-\alpha)\nmid b_{1}(s)$; set $\Gamma_{n}'=\Gamma_{n-1}+(t\partial_{t}-\alpha-n)\Gamma_{n}$, then $(t\partial_{t}-(\alpha-1)-n)^\ell b_{1}(t\partial_{t}-n)$ kills $\frac{\Gamma_{n}'}{\Gamma_{n-1}'}$ for all $n$. 

    If $\mathrm{Re}(\alpha-1)\notin(-1,0]$, repeat this until all roots moved in $(-1,0]$. 
\end{proof}

\begin{definition}
    \quad

    \begin{enumerate}
        \item  $M$ is called $Q\textbf{-specializable}$ along $Y$ if the $b$-function has all its roots in $Q$, where $Q \subset \mathbb{C}$ is a subfield.
        \item If $M$ is $Q$-specializable, then for all $\beta \in Q$, take $k = \lceil \beta \rceil$, then define $$V_{\beta} M = \pi^{-1} \left( \bigoplus_{\alpha \leq \beta} \left( \frac{V_k M}{V_{k-1} M} \right)^{\alpha} \right),$$and $$V_{<\beta} M = \pi^{-1} \left( \bigoplus_{\alpha < \beta} \left( \frac{V_k M}{V_{k-1} M} \right)^{\alpha} \right),$$where $\left( \frac{V_k M}{V_{k-1} M} \right)^{\alpha}$ is the $\alpha$-eigenspace of $\frac{V_k M}{V_{k-1} M}$ w.r.t. the $t\partial_t$-action, and $\pi: V_k M \rightarrow \frac{V_k M}{V_{k-1} M}$; then we get a $\mathbb{Q}$-indexed $V$-filtration s.t.  $$\mathrm{gr}_{\alpha}^{V} M = \frac{V^{\alpha} M}{V^{<\alpha} M} = \left( \frac{V_k M}{V_{k-1} M} \right)^{\alpha}$$ with $t\partial_t - \alpha$ acts on $\mathrm{gr}_{\alpha}^{V} M$ nilpotently. 
    \end{enumerate}
\end{definition}
\begin{example}
    Let $X=\mathbb{A}^{1}_{\mathbb{C}}$ and $M = \mathbb{C}\left[ t, \frac{1}{t} \right]e^{\frac{1}{t}}$ an irregular holonomic $\mathcal{D}$-mod.

    Consider $V$-filtration along $t = 0$. The $V$-filtration on $M$ exists by Theorem \ref{Theorem 8.7}$$\partial_t e^{\frac{1}{t}} = -\frac{1}{t^2} e^{\frac{1}{t}}\text{ which implies that }t\partial_t e^{\frac{1}{t}} = -\frac{1}{t} e^{\frac{1}{t}}(\ast)$$Take $\Omega_i M = V_i \mathcal{D}_{X} \cdot e^{\frac{1}{t}}$ for $i \in \mathbb{Z}$, then $\Omega_0 M = M$ for $(\ast)$ holds, further, we obtain that $\Omega_1 M = t \Omega_0 M = t M = M$ and $\Omega_i M = M$ for all $i\in \mathbb{Z}$. Finally, we find out that $\Omega_{i}M$ is the $V$-filtration we want with the $b$-function of $M$ is $1$. 

\end{example}

\begin{theorem}
    If M is a regular holonomic $\mathcal{D}_X$-mod, then its V-filtration is not trivial. In particular, its b-function is not $1$.
\end{theorem}

\begin{proof}[Sketch of proof]
    First, each $V_i M$ is a $V_0 \mathcal{D}_X$-module and $V_0 \mathcal{D}_X = \mathcal{D}_X^{\log}$ w.r.t. $(t = 0)$. Since $V_i M$ has no $t$-torsion for all $i < 0$, $\mathrm{ch}(V_i M)$ has no component over $(t = 0)$, for all $i < 0$. 

    Further, we obtain that $\mathrm{ch}^{\log}\left(\frac{V_i M}{V_{i-1} M}\right) = \mathrm{ch}^{\log}\left(\frac{V_i M}{t V_i M}\right)= \mathrm{ch}^{\log}(V_i M)\big|_{(t=0)}$, thus $\frac{V_i M}{V_{i-1} M} \neq 0.$ 

\end{proof}


\begin{example}\label{Example 8.11}

Let $X = \mathbb{A}_{\mathbb{C}}^n$ and $Y = \mathbb{A}_{\mathbb{C}}^n \times \mathbb{A}_{\mathbb{C}}^1$, here $\mathbb{A}^{1}_{\mathbb{C}}$ is with coordinate $t$. 
For $f \in \mathbb{C}[x_1, \cdots, x_n]$, Consider $M = \mathbb{C}[x_1, \cdots, x_n]_f[s] f^s$.

Define $t M = f M$ and then $t^{-1} M = f^{-1} M$.  Take $s = -t\partial_t$, then $\partial_t = -t^{-1} s$, the $\partial_{t}$-action is defined, and further, $M$ is a $\mathcal{D}_Y = \mathcal{D}_X \langle t, \partial_t \rangle$-module. 

Then $N := \mathbb{C}[x_1, \cdots, x_n][\partial_t] f^s = \mathbb{C}[x_1, \cdots, x_n][\partial_t] f^s= \mathbb{C}[x_1, \cdots, x_n][t^{-1} s] f^s$, and $\partial_{x_i} f^s = \partial_{x_i}(f) \cdot \frac{1}{f} s f^s= \partial_{x_i}(f) \cdot (-\partial_t) f^s$, thus $N$ is a $\mathcal{D}_Y$-submodule of $M$. Furthermore, we can check that $\mathrm{ch}(N)$ = the conormal bundle of the graph of $X$, which indicates that $N$ is a holonomic $\mathcal{D}_{Y}$-module. 

We consider $V_{i}\mathcal{D}_{Y}\cdot f^{s}$ forming a good filtration of $N$: $$V_0 \mathcal{D}_Y \cdot f^s = \mathcal{D}_X \langle t, s \rangle f^s = \mathcal{D}_X[s] f^s$$Similarly, $V_{-1} \mathcal{D}_Y \cdot f^s = \mathcal{D}_X[s] f^{s+1}$. 

By Theorem \ref{Theorem 8.4}, there exists $b(s) \in \mathbb{C}[s]$ s.t.$$
b(s) \cdot \frac{V_0 \mathcal{D}_Y \cdot f^s}{V_{-1} \mathcal{D}_Y \cdot f^s} = b(s) \cdot \frac{\mathcal{D}_X[s] f^s}{\mathcal{D}_X[s] f^{s+1}} = 0.$$which is to say there exists $P(s)\in \mathcal{D}_{X}[s]$ such that $b(s) f^s = P(s) \cdot f^{s+1}$

In fact, $N = \iota_{f,+} (\mathbb{C}[x_1, \ldots, x_n])$, the $\mathcal{D}$-module pushforward, where $$\iota_f: X \to Y = X \times \mathbb{A}^{1}_{\mathbb{C}}\atop x \mapsto (x, f(x))$$the graph embedding. So $N$ is regular holonomic, and $b(s) \neq 1$. 
\end{example}

\begin{remark}
    One can replace $\mathbb{C}[x_1, \ldots, x_n]$ by any regular holonomic $\mathcal{D}_X$-module, to get similar results in general.
\end{remark}


\chapter{Application: Nearby cycles and comparsion}

\section{Two definitions of nearby cycles}
On $\mathbb{A}_{\mathbb{C}}^1$ with coordinate $t$, given any $\alpha \in \mathbb{C}$, we define: $$K_m^{\alpha} = \bigoplus_{\ell=0}^{m-1} \mathbb{C}[t, t^{-1}] e_\ell^{\alpha}$$ with $\partial_t e_\ell^{\alpha} = \frac{1}{t}(\alpha e_\ell^{\alpha} + e_{\ell-1}^{\alpha})$ for all $\ell$, where $e_{-1}^{\alpha} = 0$. 

Then the matrix of $t\partial_t$-action on $K_m^{\alpha}$ is $J_{\alpha,m}$, the $m\times m$ Jordan block with eigenvalue $\alpha$:$$J_{\alpha,m} = \begin{bmatrix}
    \alpha & 1 & & \\
     & \alpha & \ddots & \\
     & & \ddots & 1 \\
     & & & \alpha
    \end{bmatrix}_{m \times m}$$

  \begin{remark}
    Here $e_\ell^{\alpha}$ can be understood as the formal symbol of the multivalued function "$t^{\alpha} \cdot \frac{(\log t)^\ell}{\ell!}$". And $K^\alpha_{m}$ indicates about the section not flat.  
  \end{remark}

  Then locally around a simply connected neighborhood of $\mathbb{A}_{\mathbb{C}}^1 \setminus \{0\}$, with simply connectedness, we obtain that $\mathrm{\log}t$ is well-defined. Consider $$v_{\ell} := e^{-J_{\alpha,m} \log t} \cdot e_{\ell}^{\alpha}, \quad \ell = 0, \cdots, m-1,$$we then have  $t\partial_t \left(e^{-J_{\alpha,m} \log t} \cdot e_{\ell}^{\alpha}\right) = -J_{\alpha,m} \left(e^{-J_{\alpha,m} \log t} \cdot e_{\ell}^{\alpha}\right) + J_{\alpha,m} \left(e^{-J_{\alpha,m} \log t} \cdot e_{\ell}^{\alpha}\right) = 0$, i.e. $v_{\ell}$ are $t\partial_t$-flat. 

  Since $\log t$ are multivalued, we thus have a local system: $$L_m^{\lambda}=(\bigoplus \limits_{i=0}^{m-1}\mathbb{C}\cdot v_{i}, e^{-2\pi iJ_{\alpha,m}})$$$e^{-2\pi i J_{\alpha,m}}$ is the monodromy action $T$ of the loop around $0 \in \mathbb{A}_{\mathbb{C}}^1$, and $\lambda = e^{-2\pi i \alpha}$, then
  $$ \text{DR}(K_m^{\alpha}) \stackrel{\text{quasi-isomorphic}}{\simeq} L_m^{\lambda}. $$
  Since the direct system $$ K_{m-1}^{\alpha} \rightarrow K_m^{\alpha} \rightarrow K_{m+1}^{\alpha} \rightarrow \cdots $$induces a direct system of local systems: $$ L_{m-1}^{\lambda} \rightarrow L_m^{\lambda} \rightarrow L_{m+1}^{\lambda} \rightarrow \cdots $$which is $T$-action compatible, and then we take limit. 
  
  Let $X = \text{Spec } R$ be a smooth affine algebraic variety over $\mathbb{C}$, and $f \in R$, a regular function, we have the following diagram: 
  \begin{center}
    \begin{tikzcd}
        {(f=0)\atop \text{divisor}} & X && {\mathbb{A}^1_{\mathbb{C}}} \\
        \\
        & {U=X\backslash(f=0)} && {\mathbb{A}^1_{\mathbb{C}}\backslash \{0\}}
        \arrow["i", hook, from=1-1, to=1-2]
        \arrow["f", from=1-2, to=1-4]
        \arrow["j", hook', from=3-2, to=1-2]
        \arrow["f", from=3-2, to=3-4]
        \arrow[hook', from=3-4, to=1-4]
    \end{tikzcd}
\end{center}

  
  We define $$ \Psi_{f,\lambda} = \lim_{\substack{\longrightarrow \\ m \to \infty}} i^{-1} R_{j_*} L_m^{\lambda} $$which is called the $\lambda\textbf{-nearby cycles}$ of $f$. And $T$-action on $L_m^{\lambda}$ induces $T$-action on $\Psi_{f,\lambda}$. 
  
\begin{remark}
    Here $j_{\ast}$ is the direct image functor, and it's a left exact functor, we write $R_{j_{\ast}}$ the right derived functor of $j_{\ast}$, which can be viewed as the complex of sheaves obtained by taking the injective resolution of $L_{m}^{\lambda}$ and applying $j_{\ast}$. Since $j$ maybe have nontrivial derived cohomology, we choose to use drived functor. 

    The above definition is a topological definition, which is a bit differed from Deligne's. And total cycles is defined by combining all information from nearby cycles. 
\end{remark}

On the other hand, we have $$ \mathcal{D}_{X}[s]f^{s} \longrightarrow R[\frac{1}{f}][s] \cdot f^{s} $$For closed point $\alpha \in \mathbb{A}_{\mathbb{C}}^{1} \simeq \operatorname{Spec} \mathbb{C}[s]$, we use $\mathfrak{m}_{\alpha}$ to denote its maximal ideal. We have proved in Example \ref{Example 8.11}: 

\begin{theorem}
    There exists $b(s) \in \mathbb{C}[s]$ s.t.$$b(s) \quad \text{kills} \quad \frac{\mathcal{D}_{X}[s]f^{s}}{\mathcal{D}_{X}[s]f^{s+1}}$$
\end{theorem}
 


\begin{corollary}[Bernstein-Sato]
    Let $\mathbb{C}(s)$ be the fractional field of $\mathbb{C}[s]$, and $\mathcal{D}_{X}(s) = \mathcal{D}_{X} \otimes_{\mathbb{C}[s]} \mathbb{C}(s)$; for all $\alpha \in \mathbb{A}_{\mathbb{C}}^{1}$ closed point, we obtain that:
   \begin{enumerate}
    \item $\mathcal{D}_{X}(s) f^{s+k} \simeq \mathcal{D}_{X}(s) f^{s}, $ for all $  k \in \mathbb{Z}$.
    \item There exists $k_{0} \geqslant 0$ s.t.
    \begin{itemize}
        \item $\mathcal{D}_{X}[s]_{\mathfrak{m}_{\alpha}} f^{s-k} = \mathcal{D}_{X}[s]_{\mathfrak{m}_{\alpha}} f^{s-k_{0}}, $ for all $ k > k_{0}$
        \item $\mathcal{D}_{X}[s]_{\mathfrak{m}_{\alpha}} f^{s+k} = \mathcal{D}_{X}[s]_{\mathfrak{m}_{\alpha}} f^{s+k_{0}},$ for all $  k > k_{0}$
    \end{itemize}
   \end{enumerate} 
  
\end{corollary}
\begin{remark}
    Noted that here  $\mathcal{D}_{X}[s] f^{s} = R[\frac{1}{f}][s] f^{s},\ \mathcal{D}_{X}[s]_{\mathfrak{m}_{\alpha}} f^{s} = R[\frac{1}{f}][s]_{m_{\alpha}} f^{s}$. 
\end{remark}
\begin{proof}
    $1)$ Since $\mathbb{C}(s)$ can be view as the localization at $(0)$ of $\mathbb{C}[s]$, we have $\mathbb{C}(s)$ is flat over $\mathbb{C}[s]$, thus $\frac{\mathcal{D}_{X}[s]f^{s}}{\mathcal{D}_{X}[s]f^{s+1}}\otimes \mathbb{C}(s)=\frac{\mathcal{D}_{X}(s)f^{s}}{\mathcal{D}_{X}(s)f^{s+1}}$.  And $b(s+k)$ is invertible in $\mathbb{C}(s)$, thus that $b(s+k)$ killing $\frac{\mathcal{D}_{X}(s)f^{s}}{\mathcal{D}_{X}(s)f^{s+1}}$ implies that $\frac{\mathcal{D}_X(s) f^{s+k}}{\mathcal{D}_X(s) f^{s+k-1}} = 0$. 

    $2)$  Know that there exists $k_0 \ge 0$ s.t. $b(s \pm k) \notin \mathfrak{m}_\alpha$ for all $k > k_0$, so $b(s \pm k)$ is invertible in $\mathbb{C}[s]_{\mathfrak{m}_\alpha}$ for all $k > k_0$, similarly, we get $\frac{\mathcal{D}_X[s]_{\mathfrak{m}_\alpha} f^{s \pm k}}{\mathcal{D}_X[s]_{\mathfrak{m}_\alpha} f^{s \pm k + 1}} = 0$. 
\end{proof}

\begin{definition}
    For all $\alpha \in \mathbb{C}$:
    \begin{enumerate}
        \item $\mathcal{D}_X[s]_{m_\alpha} f^{s+k},  k \gg 0$ is stabled, called the $\textbf{minimal extension}$ of $R[\frac{1}{f}][s]_{m_\alpha} f^s$. 
        \item $\Psi_{f,\alpha} := \frac{D_X[s]_{m_\alpha} f^{s-k}}{D_X[s]_{m_\alpha} f^{s+k}}$ is called the $\alpha\textbf{-nearby cycles}$ of $R$ along $f$. 
        \item $\Lambda := \mathbb{Z} + \{\text{roots of } b_{f}(s)\}$. 
    \end{enumerate}
\end{definition}

\begin{remark}
    In $\mathcal{D}$-module theory, an extension is usually about given a module over an open subset $U$, and extending it to the whole space $X$ to get a $\mathcal{D}_{X}$-module, and we shall pay attention on how it extend around $f=0$. Take $U=\mathrm{Spec}R_{f}$, if the module over $U=X\backslash(f=0)$ can be viewed as $R$-module, then it's an extension, and we call $R_{f}[s]_{\mathfrak{m}_{\alpha}}f^{s}$  $\textbf{maximal extension}$. 
\end{remark}

\begin{proposition}
    \quad

    \begin{enumerate}
        \item $s - \alpha$ acts on $\Psi_{f,\alpha}$ nilpotently
        \item $\Psi_{f,\alpha} \cong \Psi_{f,\alpha+1}$
        \item $\Psi_{f,\alpha} \neq 0$ if and only if $\alpha \in \Lambda$
        \item $\Psi_{f,\alpha}$ is a regular holonomic $D_X$-module. 
    \end{enumerate}
\end{proposition}
\begin{proof}
    $1)$  Since we have a chain $\mathcal{D}_{X}[s]_{\mathfrak{m}_{\alpha}}f^{s+k_{0}}\subseteq\mathcal{D}_{X}[s]_{\mathfrak{m}_{\alpha}}f^{s+k_{0}-1}\subseteq\dots\subseteq \mathcal{D}_{X}[s]_{\mathfrak{m}_{\alpha}}f^{s-k_{0}}$, we can write $\frac{\mathcal{D}_{X}[s]_{\mathfrak{m}_{\alpha}}f^{s-k_{0}}}{\mathcal{D}_{X}[s]_{\mathfrak{m}_{\alpha}}f^{s+k_{0}}}=\frac{\mathcal{D}_{X}[s]_{\mathfrak{m}_{\alpha}}f^{s-k_{0}}}{\mathcal{D}_{X}[s]_{\mathfrak{m}_{\alpha}}f^{s-k_{0}+1}}\cdots\frac{\mathcal{D}_{X}[s]_{\mathfrak{m}_{\alpha}}f^{s}}{\mathcal{D}_{X}[s]_{\mathfrak{m}_{\alpha}}f^{s+1}}\cdots\frac{\mathcal{D}_{X}[s]_{\mathfrak{m}_{\alpha}}f^{s+k_{0}-1}}{\mathcal{D}_{X}[s]_{\mathfrak{m}_{\alpha}}f^{s+k_{0}}}$ to show that $\frac{\mathcal{D}_{X}[s]_{\mathfrak{m}_{\alpha}}f^{s-k_{0}}}{\mathcal{D}_{X}[s]_{\mathfrak{m}_{\alpha}}f^{s+k_{0}}}$ is determined by those intermediate quotient modules. And for $\frac{\mathcal{D}_{X}[s]_{\mathfrak{m}_{\alpha}}f^{s+i}}{\mathcal{D}_{X}[s]_{\mathfrak{m}_{\alpha}}f^{s+i+1}}$, know that $b_{i}(s)$ exists, and since only $s-\alpha$ in not invertible, $b_{i}(s)$ must take the form of $m_{i}(s-\alpha)^{n_{i}}$. We take the product of all $b_{i}(s)$, or just take $n\gg \sum n_{i}$, we will have $(s-\alpha)^{n_{i}}$ annihilate $\frac{\mathcal{D}_{X}[s]_{\mathfrak{m}_{\alpha}}f^{s-k_{0}}}{\mathcal{D}_{X}[s]_{\mathfrak{m}_{\alpha}}f^{s+k_{0}}}$. 
	
	$2)$ Substitute $s$ by $s+1$, all shifts, nothing changes. 
	
	$3)$ Since $\alpha \not\in\Lambda$, $\alpha$ is not a root of $b(s+k)$, then $b(s+k)$ is a unit, thus $\Psi_{f,\alpha}=0$. 
	If $\Psi_{f,\alpha}\neq 0$, there exists $i$ s.t. $\frac{\mathcal{D}_{X}[s]_{\mathfrak{m}_{\alpha}}f^{s+i}}{\mathcal{D}_{X}[s]_{\mathfrak{m}_{\alpha}}f^{s+i+1}}\neq 0$, and then $\alpha$ is the root of $b(s+i)$ s.t. $b(s+i)$ is not invertible. Thus $\alpha\in\Lambda$. 
	
	$4)$ $M_{\alpha} = \mathbb{C}[t, \frac{1}{t}] \cdot t^{\alpha}$ is regular, holonomic $\mathcal{D}_{\mathbb{A}^{1}_{\mathbb{C}}}$-module, consider $$f:X\rightarrow \mathbb{A}^{1}_{\mathbb{C}}\atop \alpha\mapsto t=f(\alpha)$$which is induced by $\mathbb{C}[t]\rightarrow R$ with $t\rightarrow f$. 
	We admit the fact that pull-back functor preserves regular holonomicity, then $f^{*} M_{\alpha}:=M_{\alpha}\otimes_{\mathbb{C}[t]} R\cong R_{f}\cdot f^\alpha= R[\frac{1}{f}] f^{\alpha}$ makes $R_{f}\cdot f^\alpha$ a regular holonomic $\mathcal{D}_{X}$-module. 
	
	With $2)$, we can assume $\alpha\neq 0$, otherwise, we just shifts $\alpha$ to $\alpha+1$. 
	We have a SES:$$0 \to \frac{R_f[s]_{m_{\alpha}} f^{s}}{(s-\alpha)^1 R_f[s]_{m_{\alpha}} f^{s}}  \to \frac{R_f[s]_{m_{\alpha}} f^{s}}{(s-\alpha)^2 R_f[s]_{m_{\alpha}} f^{s}} \xrightarrow{s-\alpha} \frac{R_f[s]_{m_{\alpha}} f^{s}}{(s-\alpha)^1 R_f[s]_{m_{\alpha}} f^{s}}  \to 0$$With $\frac{R_f[s]_{m_{\alpha}} f^{s}}{(s-\alpha)^1 R_f[s]_{m_{\alpha}} f^{s}}=R_{f} \cdot f^{\alpha}$ regular holonomic. 
	
	Since regular holonomic $\mathcal{D}_{X}$-module forms an abelian category,  we have $\frac{R_{f}[s]_{\mathfrak{m}_{\alpha}}f^{s}}{(s-\alpha)^{2}R_{f}[s]_{\mathfrak{m}_{\alpha}}f^{s}}$ is regular holonomic. 
	
	And further, we obtain that: $$0 \to \frac{R_f[s]_{m_{\alpha}} f^{s}}{(s-\alpha)^1 R_f[s]_{m_{\alpha}} f^{s}}  \to \frac{R_f[s]_{m_{\alpha}} f^{s}}{(s-\alpha)^{N+1} R_f[s]_{m_{\alpha}} f^{s}} \to \frac{R_f[s]_{m_{\alpha}} f^{s}}{(s-\alpha)^N R_f[s]_{m_{\alpha}} f^{s}}  \to 0$$By induction, $\frac{R_f[s]_{m_{\alpha}} f^{s}}{(s-\alpha)^{N} R_f[s]_{m_{\alpha}} f^{s}}$ is regular holonomic for all $N$. 
	
	By $1)$, $\Psi_{f,\alpha}$ is a quotient of $\frac{R_f[s]_{m_{\alpha}} f^{s}}{(s-\alpha)^N R_f[s]_{m_{\alpha}} f^{s}}$, so $\Psi_{f,\alpha}$ is also regular holonomic over $\mathcal{D}_{X}$. 
\end{proof}
\section{Comparison theorem}
Our next goal is to prove the following theorem: 
\begin{theorem}[Comparison] \label{Theorem 9.5}
    For all $\alpha \in \mathbb{C}$, $\lambda = e^{2\pi i\alpha}$, $$\mathrm{DR}(\Psi_{f,-\alpha}) \simeq i_* \Psi_{f,\lambda}[-1]$$s.t. $e^{2\pi is} = \lambda^{-1}\cdot e^{2\pi i(s+\alpha)}$ is the corresponding $T$-action. 
\end{theorem}
\begin{remark}
    with $(s-\alpha)$ acts on $\Psi_{f,\alpha}$ nilpotent, $e^{2\pi i(s+\alpha)}$ will be a well-defined operator for $\Psi_{f,-\alpha}$. 
\end{remark}


Define $N_m^{\alpha,k} = \bigoplus_{\ell=1}^m \mathcal{D}_X[s] f^{s-k} \cdot e_\ell^{\alpha}$, know that $s=-t\partial_{t}$,  $s \cdot (\eta \cdot e_\ell^{\alpha}) = (s\eta) \cdot e_\ell^{\alpha}+\eta(-\alpha e^{\alpha}_{\ell}-e^{\alpha}_{\ell-1})=(s-\alpha)\eta \cdot e^{\alpha}_{\ell}-\eta \cdot e^{\alpha}_{\ell-1}$, so $N_m^{\alpha,k}$ is also an $\mathcal{D}_X[s]$-mod. 

We define $b_f(s)$ to be the monic polynomial of least degree among $b(s)$ s.t. $$b(s) \cdot \frac{\mathcal{D}_X[s] f^s}{\mathcal{D}_X[s] f^{s+1}} = 0$$called the $\textbf{b-function} \text{ or } \textbf{Bernstein-Sato polynomial of }f$.

\begin{lemma}\label{lemma 9.6}
    \quad

    \begin{enumerate}
        \item $\mathrm{DR}\left(R_f[s] f^s \otimes^L_{\mathbb{C}(s)} \mathbb{C}_\alpha\right) \simeq \mathrm{DR}\left(R_f[s] f^{s+\alpha} \otimes^L_{\mathbb{C}(s)} \mathbb{C}_\alpha\right)\simeq R j_* (f^{-1} L_1^\lambda)$
        \item $\mathrm{DR}\left(\mathcal{D}_X[s] f^{s+k} \otimes^{L}_{\mathbb{C}(s)} \mathbb{C}_\alpha\right) \simeq \mathrm{DR}\left(\mathcal{D}_X[s] f^{s+k+\alpha} \otimes^{L}_{\mathbb{C}(s)} \mathbb{C}_\alpha\right)\simeq j_! (f^{-1} L_1^\lambda), \quad k \gg |\alpha|$
    \end{enumerate}
    where $\mathbb{C}_\alpha$ is the residue field of $\alpha \in \mathbb{A}_{\mathbb{C}}^1$, and $\lambda = e^{2\pi i \alpha}$.
\end{lemma}
\begin{proof}
    Take $0\rightarrow \mathbb{C}[s]\xrightarrow{s-\alpha} \mathbb{C}[\alpha]\rightarrow \frac{\mathbb{C}[s]}{s-\alpha}=\mathbb{C}_{\alpha}$ free resolution, then we obtain that $R_f[s]f^s \otimes^L_{\mathbb{C}(s)} \mathbb{C}_\alpha$ $$\stackrel{\text{q.i}}{\simeq}[ R_f[s]f^s \xrightarrow{s-\alpha} R_f[s]f^s ]\stackrel{\text{q.i}}{\simeq}[R_{f}[s]f^{s+\alpha}\xrightarrow{s} R_f[s]f^{s+\alpha}]\simeq R_f f^\alpha.$$with $U\xrightarrow{f}\mathbb{A}^{1}_{\mathbb{C}}\backslash\{0\}$, $\mathrm{DR}(R_f \cdot f^\alpha)|_U \simeq f^{-1} L_1^\lambda$, using log resolution of $(X,f=0)$, one can prove $\mathrm{DR}(R_f \cdot f^\alpha) \simeq R j_*(f^{-1} L_1^\lambda)$. 
Using resolution of singularities and duality, one can prove 2.

\end{proof}
\begin{lemma} \label{lemma 9.7}
    $$(b_f(s-k-\alpha))^m \text{ kills } N_m^{\alpha,k}/N_m^{\alpha,k-1}$$
\end{lemma}
\begin{proof}
    By construction, we have a SES: $$0\longrightarrow N_{m-1}^{\alpha,k} \longrightarrow N_m^{\alpha,k} \longrightarrow Q \longrightarrow 0$$with $Q=\mathcal{D}_{X}[s]f^{s-k}\bar{e}_{m-1}^{\alpha}$, with $s\cdot \bar{e}^\alpha_{m-1}=(s-\alpha)\overline{e}^\alpha_{m-1}$, we choose $s-\alpha$ to replace $s$, then  $Q\simeq \mathcal{D}_{X}[s]f^{s-k-\alpha}$. 

    Since $b_{f}(s-k-\alpha)$ kills $\mathcal{D}_{X}[s]f^{s-k-\alpha}$, and by induction, we have $(b_f(s-k-\alpha))^{m -1}$ kill $N_{m-1}^{\alpha,k}$, done. 
\end{proof}
\begin{proposition}
    There exists $k_0 > 0$ s.t. for all $k \geq k_0$, we obtain that:
    \begin{enumerate}
        \item $\mathrm{DR}(N_m^{-\alpha,k} \xrightarrow{s} N_{m}^{-\alpha,k}) \simeq R j_*\left(f^{-1} L_m^{\lambda}\right)$; 
        \item $\mathrm{DR}(N_m^{-\alpha,-k} \xrightarrow{s} N_{m}^{-\alpha,-k})\simeq j_!\left(f^{-1} L_m^{\lambda}\right)$. 
    \end{enumerate} 
\end{proposition}
    \begin{proof}
        Let $\underline{\mathfrak{m}}_0$ be the maximal ideal of $0 \in \mathbb{A}_{\mathbb{C}}^1$, then $$[N_{m}^{-\alpha,k} \xrightarrow{s} N_{m}^{-\alpha,k}]\stackrel{\text{q.i}}{\simeq}N^{-\alpha,k} \otimes^{L}_{\mathbb{C}[s]} \mathbb{C}_{0} \stackrel{\text{q.i}}{\simeq} [N^{-\alpha,k}_{m,\underline{\mathfrak{m}}_{0}} \xrightarrow{s}N^{-\alpha,k}_{m,\underline{\mathfrak{m}}_{0}}]\text{ localization}\atop \stackrel{\text{q.i}}{\simeq}[\bigoplus \limits_{\ell=1}^{m-1} R_{f}[s]f^{s} \cdot e_{\ell}^{-\alpha} \xrightarrow{s} \bigoplus \limits_{\ell=1}^{m-1}R_{f}[s]f^{s} \cdot e_{\ell}^{-\alpha}]\text{ by Lemma \ref{lemma 9.7}}.$$Using Lemma \ref{lemma 9.6} (1) for $\mathrm{DR}(R_{f}[s]f^{s+\alpha}\otimes_{\mathbb{C}[s]}\mathbb{C}_{0})$, by induction, Part (1) is proved. And part (2) is similar.
    \end{proof}

\begin{lemma}
    Let $W$ be a $\mathbb{C}$-vector space, and let $\varphi$ be a $\mathbb{C}$-linear operator on $W$ admitting a minimal polynomial.
Set $W_{\infty} = \bigoplus_{k=0}^{\infty} W \otimes e_k$ and define $$\varphi_{\infty}(w \otimes e_i) = (\varphi - \alpha) \cdot w \otimes e_i - w \otimes e_{i-1} \quad \text{for all }  w \in W, \; (e_{-1} = 0).$$Then $\varphi_{\infty}$ is surjective and $\ker(\varphi_{\infty}) = W_{\alpha}$ is the generalized $\alpha$-eigenspace.
\end{lemma}
\begin{proof}
    Define a map
    \begin{equation}
        W_{\alpha} \longrightarrow \ker(\varphi_{\infty})\atop w \longmapsto \sum\limits_{i \geqslant 0} (\varphi - \alpha)^i w \otimes e_{i}.
        \tag{$\ast$}
        \label{eq}
    \end{equation}
    This map is an isomorphism.

    Then we prove surjectivity: If $w\in W_{\alpha}$, then $\varphi_{\infty}( -\sum\limits_{i>j} (\varphi-\alpha)^{i-j-1}w\otimes e_{i})=w\otimes e_{j}$; If $w\in W_{\alpha}^{\perp}$, then $\varphi_{\infty}( \sum\limits_{i=1}^{j}(\varphi-\alpha)^{-i}w\otimes e_{j-i+1} )=w\otimes e_{j}$. 
\end{proof}
\begin{corollary}
    For $k \gg 0$, $$
\Psi_{f,\alpha}[1] \stackrel{\text{q.i}}{\simeq}\varinjlim_m \left( \frac{N_m^{\alpha,k}}{N_m^{\alpha,-k}} \xrightarrow{s} \frac{N_m^{\alpha,k}}{N_m^{\alpha,-k}} \right).
$$
\end{corollary}

\begin{proof}[Proof of  (\ref{Theorem 9.5})]
    $$
\mathrm{DR}(\Psi_{f,-\alpha}[1]) \cong \mathrm{DR}\left( \varinjlim_m  \left( \frac{N_m^{-\alpha,k}}{N_m^{-\alpha,-k}} \xrightarrow{s} \frac{N_m^{-\alpha,k}}{N_m^{-\alpha,-k}} \right) \right)$$
$$ \cong \varinjlim_m  \mathrm{DR}\left( \frac{N_m^{-\alpha,k}}{N_m^{-\alpha,-k}} \xrightarrow{s} \frac{N_m^{-\alpha,k}}{N_m^{-\alpha,-k}} \right)$$
$$\simeq \varinjlim_m\ \text{cone}\left( DR\left(N_{m}^{-\alpha,-k} \xrightarrow{s} N_m^{-\alpha,-k}\right) \rightarrow DR\left(N_m^{-\alpha,k} \xrightarrow{s} N_m^{-\alpha,k}\right)\right)$$
$$\simeq\varinjlim_m\ \text{cone}\left(j_! f^{-1} L_m^{\lambda} \rightarrow R j_* f^{-1} L_m^{\lambda}\right)\simeq \Psi_{f,\lambda}$$

By $\eqref{eq}$, we consider $$   \Psi_{f,-\alpha} \longrightarrow \ker(-s_{\infty})\atop \eta \longmapsto \sum\limits_{i \geqslant 0} (-s-\alpha)^i \eta\otimes  e_{i}.$$

Then $s+\alpha = -J_{0,\infty}$, and $e^{2\pi i(s+\alpha)} = e^{-2\pi i J_{0,\infty}}$. Furthermore, we obtain that $$e^{-2\pi i\alpha} \cdot e^{2\pi i(s+\alpha)} = e^{-2\pi i J_{\alpha,\infty}} = T$$
\end{proof}
\end{document}
